% =====================================================================
%  Paper C of the odd-zeta series.
%  Working title:
%    p-adic Apery limits: rate laws, valuation strata, and closed forms
%    via the p-adic Gamma function
%
%  DRAFT.  Author placeholder on line \author{...} below.
%  Evidence classes are printed in the statements; see Section 1.5.
% =====================================================================
\documentclass[11pt,reqno]{article}

\usepackage[T1]{fontenc}
\usepackage[utf8]{inputenc}
\usepackage[margin=1.15in]{geometry}
\usepackage{amsmath,amssymb,amsthm,mathtools}
\usepackage{booktabs}
\usepackage{array}
\usepackage{longtable}
\usepackage{enumitem}
\usepackage{microtype}
\usepackage{xcolor}
\usepackage[colorlinks=true,linkcolor=blue!55!black,citecolor=blue!55!black,%
            urlcolor=blue!55!black]{hyperref}

% ---------------------------------------------------------------- macros
\newcommand{\Qp}{\mathbb{Q}_p}
\newcommand{\Zp}{\mathbb{Z}_p}
\newcommand{\QQ}{\mathbb{Q}}
\newcommand{\ZZ}{\mathbb{Z}}
\newcommand{\CC}{\mathbb{C}}
\newcommand{\NN}{\mathbb{N}}
\newcommand{\Gp}{\Gamma_p}
\newcommand{\Ph}{\widehat{P}}
\newcommand{\Ih}{\widehat{I}}
\newcommand{\Ip}{I'}
\newcommand{\Lah}{\Lambda^{\Ph}}
\newcommand{\Lap}{\Lambda^{P}}
\newcommand{\Atil}{\widetilde{A}}
\newcommand{\Btil}{\widetilde{B}}
\newcommand{\vp}{v_p}
\DeclareMathOperator{\ord}{ord}
\DeclareMathOperator{\lcm}{lcm}
\DeclareMathOperator{\cond}{cond}
\DeclareMathOperator{\disc}{disc}
\newcommand{\bin}[2]{\binom{#1}{#2}}
\newcommand{\floor}[1]{\lfloor #1 \rfloor}

% evidence tags
\newcommand{\PROVED}{{\normalfont\textsc{[proved]}}}
\newcommand{\VER}[1]{{\normalfont\textsc{[verified:} #1\textsc{]}}}
\newcommand{\RECALLED}{{\normalfont\textsc{[recalled--unverified]}}}
\newcommand{\FIT}{{\normalfont\textsc{[conjectural fit]}}}
\newcommand{\OPENTAG}{{\normalfont\textsc{[open]}}}

% ---------------------------------------------------------------- theorems
\theoremstyle{plain}
\newtheorem{theorem}{Theorem}[section]
\newtheorem{proposition}[theorem]{Proposition}
\newtheorem{lemma}[theorem]{Lemma}
\newtheorem{corollary}[theorem]{Corollary}
\newtheorem{conjecture}[theorem]{Conjecture}

\theoremstyle{definition}
\newtheorem{finding}[theorem]{Finding}
\newtheorem{definition}[theorem]{Definition}
\newtheorem{example}[theorem]{Example}
\newtheorem{problem}[theorem]{Problem}

\theoremstyle{remark}
\newtheorem{remark}[theorem]{Remark}
\newtheorem{question}[theorem]{Question}

\numberwithin{equation}{section}

% ---------------------------------------------------------------- title
\title{\bfseries $p$-adic Ap\'ery limits:\\ rate laws, valuation strata,\\ and
closed forms via the $p$-adic Gamma function}

% ---------------------------------------------------------------------
% AUTHOR PLACEHOLDER -- to be filled in.
% ---------------------------------------------------------------------
\author{[AUTHOR NAME(S) AND AFFILIATION TO BE SUPPLIED]}
\date{Draft, \today}

\begin{document}
\maketitle

\begin{abstract}
Let $(A(n),B(n))$ be an Ap\'ery-like pair: $A$ one of the fifteen sporadic
binomial sums, $B$ the second solution of its recurrence.  Classically
$B(n)/A(n)$ converges to the \emph{Ap\'ery limit}, a rational multiple of
$\zeta(w)$ or of $L(\chi,w)$, where $w$ is the weight and $\chi$ a quadratic
character attached to the sequence.  We study the $p$-adic analogue.  Along a
$p$-power tower $n_s=ap^s$ the normalised ratios $\chi(p)^{s}p^{ws}B(n_s)/A(n_s)$
converge in $\Qp$; we call the limit $\Lambda_a$ a \emph{$p$-adic Ap\'ery
limit}.  Convergence is driven by the digit-descent congruences
$\vp\bigl(p^{w}B_nA_q-\chi(p)B_qA_n\bigr)\ge w$, $q=\floor{n/p}$, of the
companion Paper~A.

We report four results.  \textbf{(1) A rate law.}  The number of new certified
$p$-adic digits per tower level is $\min\bigl(3,\,r(\chi,w)\bigr)$, where the
$3$ is the Jacobsthal--Kazandzidis exponent governing $\bin{Ap^t}{Bp^t}$ and
$r(\chi,w)$ is the rate of the $\chi$-twisted harmonic tower.  The rate is
therefore \emph{not} the motivic weight, correcting a natural guess; for weight
$5$ the only correct multi-digit statement is the scaled congruence
$p^{5s}P_{ap^s}Q_{ap^{s-1}}\equiv p^{5(s-1)}P_{ap^{s-1}}Q_{ap^{s}}
\pmod{p^{3s}}$, of depth $3s$.  \textbf{(2) Valuation strata and a vanishing
defect.}  For the Brown--Zudilin $\zeta(5)$ family, $\vp(\Ph_n/Q_n)=-3L+O(1)$
and $\vp(P_n/Q_n)=-5L+O(1)$ with $L=\floor{\log_p n}$ and the $O(1)$ bounded by
$2$; consequently the $p$-adic purity defect vanishes identically,
$c_p=\lim_n \Ih_n/\Ip_n=0$, in sharp contrast with the archimedean defect
$c=-3/\pi^2$.  \textbf{(3) An exclusion theorem.}  No tower limit of the
classical Ap\'ery pair or of the Brown--Zudilin pair is a $\QQ$-affine function
of $\zeta_p(3)$ or $\zeta_p(5)$ of height $\lesssim 10^4$; structurally, the
$p$-adic linear forms in $\zeta_p(3)$, $\zeta_p(5)$ built from these pairs
\emph{diverge polynomially}, so no sharpening can rescue them.  These pairs do
not $p$-adically approximate the Kubota--Leopoldt zeta values at all.
%% FIXED (MAJOR M-1, evidence-class): "We prove" was attached to the one assembly that is only
%% FIXED [VERIFIED] --- Finding 6.6, \VER{p=5,7,11,13, block by block}; LAMBDA_HUNT:185 labels it
%% FIXED "[VERIFIED order by order]", never [PROVED]. Only h_p (Thm 6.1) and c_p (Thm 6.3) are
%% FIXED proved, and the abstract now says exactly that.
\textbf{(4) An identification that explains (3).}  We exhibit, and verify block
by block, a decomposition
$\Lambda_a=h_p(a)+(\text{blocks})/\Atil_a$ with
$h_p(a)=\sum_{x\in\ZZ[1/p],\,0<x\le a}x^{-3}$, every block a polynomial in the
Kazandzidis limits $c_p(A,B)=\lim_t\bin{Ap^t}{Bp^t}$; and we \emph{prove} the two
ingredients of that decomposition --- the closed form of $h_p$, and $c_p(A,B)$ in
closed form: $-1$ times a $p$-power times a ratio of Morita $\Gp$-values times $\exp$ of an
explicit $\QQ$-series in \emph{all} odd $p$-adic zeta values.  Thus $\Lambda$ is
a \emph{value} of the Beukers--Vlasenko generating function
$\Gp(x)e^{-\Gp'(0)x}=\exp\bigl(-\sum_{m\ge2}\zeta_p(m)x^m/m\bigr)$ at $p$-power
arguments, whereas the Beukers--Vlasenko Frobenius structure constants are its
Taylor coefficients at $0$ --- and for the order-$3$ Ap\'ery operator those
slots are literally zero.

Finally, three of the fifteen families have no archimedean Ap\'ery limit at all
--- their characteristic roots are complex conjugates of equal modulus --- and
yet they do have $p$-adic tower limits.
\end{abstract}

\tableofcontents

% =====================================================================
\section{Introduction}\label{sec:intro}
% =====================================================================

\subsection{Ap\'ery limits, classically}\label{ss:classical}

Ap\'ery's irrationality proof for $\zeta(3)$ \cite{Apery1979,Fischler2003}
produces two solutions of one three-term recurrence,
\begin{equation}\label{eq:apery-rec}
  (n+1)^3u_{n+1}=(34n^3+51n^2+27n+5)\,u_n-n^3u_{n-1},
\end{equation}
namely the integers
\[
  a_n=\sum_{k}\bin{n}{k}^2\bin{n+k}{k}^2
\]
and the second solution $b_n$ normalised by $b_0=0$, $b_1=6$; the pair satisfies
$b_n/a_n\to\zeta(3)$.  The number $\zeta(3)$ appears here not as an integral but
as a \emph{limit of a ratio of two solutions of a difference equation}.  This is
the phenomenon that has come to be called an \emph{Ap\'ery limit}.

It is not isolated.  Zagier's search for integral solutions of the
$(n+1)^2$-recurrence \cite{Zagier2009}, the Almkvist--Zudilin list at
$(n+1)^3$, and Cooper's three additional cases \cite{Cooper2012}
together give fifteen \emph{sporadic} Ap\'ery-like pairs.  We use throughout the
labelling and data of Malik--Straub \cite{MalikStraub} and Straub
\cite{Straub2023}, reproduced in Table~\ref{tab:families}.  Each pair has a
weight $w\in\{2,3\}$ --- the minimal exponent with $d_n^{\,w}B(n)\in\ZZ$,
$d_n=\lcm(1,\dots,n)$ --- and an Ap\'ery limit which is a rational multiple of
$\zeta(w)$ or of $L(\chi_D,w)$ for a quadratic character $\chi_D$.  The
character of the limit turns out to be an invariant of the whole $p$-adic story
as well; see Finding~\ref{find:twist}.

Beyond the sporadic list, the same shape occurs in the Brown--Zudilin cellular
approximations to $\zeta(5)$ \cite{BrownZudilin2022}, whose conceptual setting
--- Ap\'ery-type proofs as cell integrals on moduli spaces, and the structural
reason $\zeta(5)$ resists --- is Brown's \cite{Brown2014}.  There one has a rank-$3$
system with rows $Q_n$, $\Ph_n$, $P_n$ and \emph{two} Ap\'ery limits:
$\Ph_n/Q_n\to\zeta(3)$ and $P_n/Q_n\to\zeta(5)$.  This is our weight-$5$ test
case throughout.

\subsection{What is known $p$-adically}\label{ss:known}

Almost nothing, and in particular nothing named.  The literature contains an
extensive study of the \emph{convergence} of the towers we are about to use ---
the supercongruences $A(mp^{r})\equiv A(mp^{r-1}) \pmod{p^{3r}}$ of Beukers,
Coster and Ishikawa and their descendants \cite{Supercong}, and the
Dwork-crystal framework of Beukers--Vlasenko \cite{BVdwork3} which explains the
$A$-row congruences as the unit-root entry of a Frobenius matrix --- but the
\emph{limits} of the corresponding second-solution towers do not appear to have
been studied at all.  Our own searches turned up no name, no notation and no
closed form for them; see \S\ref{ss:BVposition} for what we did and did not
find.

There are two nearby bodies of work, and it is important to say precisely how
our object differs from each.

\begin{itemize}[leftmargin=1.5em]
\item \emph{$p$-adic zeta values and irrationality.}  Calegari
  \cite{Calegari2005} proved $\zeta_p(3)\notin\QQ$ for $p=2,3$, and
  Lai--Sprang--Zudilin \cite{LSZ2025} proved $\zeta_2(5)\notin\QQ$; the
  asymptotic $p$-adic Ball--Rivoal theory is developed in
  \cite{LaiSprang2023}.  A recurring expectation --- ``one ratio, two
  completions'' --- is that the Ap\'ery pair should approximate $\zeta_p(3)$ as
  it approximates $\zeta(3)$.  Section~\ref{sec:exclusion} shows this is
  \emph{false}, and false for a structural reason.  Consistently, Calegari's
  proof does not use the Ap\'ery pair at all.
\item \emph{Frobenius constants.}  There is a genuine terminological collision
  here.  The Bloch--Vlasenko / Roy--Vlasenko \emph{Frobenius constants} $\kappa_n$
  \cite{RoyVlasenko,Kerr} are archimedean periods.  The Beukers--Vlasenko
  \emph{Frobenius structure constants} $\alpha_j$ \cite{BV2025} are $p$-adic,
  and are Taylor coefficients of $\Gp(x)e^{-\Gp'(0)x}$ at $x=0$.  Our
  $\Lambda_a$ is the $p$-adic analogue of neither: it is the analogue of the
  \emph{Ap\'ery limit itself}, i.e.\ of a $\kappa$-type period.  It happens to
  be built from the same generating function as the $\alpha_j$, but by
  \emph{evaluation} at $p$-power arguments rather than by Taylor expansion at
  $0$.  See \S\ref{ss:BVposition}.
\end{itemize}

\subsection{Results}\label{ss:results}

Fix an odd prime $p\ge5$ and a tower base $a$ with $p\nmid a$, and set
$n_s=ap^s$.

\paragraph{(A) Existence and the rate law (\S\ref{sec:engine}--\S\ref{sec:rate}).}
The digit-descent congruences of Paper~A \cite{PaperA-inprep},
\[
  (\mathrm{LB}_w^\chi)\qquad
  \vp\bigl(p^{w}B_nA_q-\chi(p)B_qA_n\bigr)\ \ge\ w,\qquad q=\floor{n/p},\qquad p\ge5,
\]
make the sequence $L_s:=\chi(p)^{s}p^{ws}B(n_s)/A(n_s)$ Cauchy, so that
\[
  \Lambda_a\ :=\ \lim_{s\to\infty}\chi(p)^{s}\,p^{ws}\,\frac{B(ap^s)}{A(ap^s)}\ \in\ \Qp
\]
exists (Theorem~\ref{thm:exists}).  The limit is an intrinsic invariant of the
recurrence and the tower: it is unchanged by $B\mapsto B+cA$, scales by
$\lambda$ under $B\mapsto\lambda B$, and satisfies
$\Lambda_{p^ka}=\chi(p)^{k}p^{-wk}\Lambda_a$
(Proposition~\ref{prop:invariance}).  The
observed convergence rate is
\[
  \boxed{\ \text{digits per level}\ =\ \min\bigl(3,\ r(\chi,w)\bigr)\ }
\]
where $r(\chi,w)$ is the rate of the $\chi$-twisted harmonic tower
(Table~\ref{tab:harmonicrate}): $r(\chi,w)=w$ for $\chi\ne1$, and $r(1,w)=w+1$,
improved to $w+2$ for odd $w$ because $\zeta_p(w+1)=0$.  The $3$ is
\emph{Jacobsthal--Kazandzidis}, not the motivic weight
(Finding~\ref{find:ratelaw}).  A first reading of the weight-$3$ and weight-$5$
data --- ``the rate is always exactly $3$'' --- is an artefact of testing only
untwisted pairs; the $\chi$-twisted weight-$2$ families converge at rate $2$.

\paragraph{(B) The weight-$5$ correction (\S\ref{ss:weight5}).}
For the Brown--Zudilin $P$-row the normalising power is $p^{5s}$ --- $\vp(p^{5s}P_{n_s}/Q_{n_s})$
is constant in $s$ --- but the single-step ratio descent $p^5f(n_s)\equiv
f(n_{s-1})$ has depth $5-2s$, which goes negative at $s\ge3$.  The surviving
multi-digit law is the scaled one,
\[
  p^{5s}\,P_{ap^{s}}\,Q_{ap^{s-1}}\ \equiv\ p^{5(s-1)}\,P_{ap^{s-1}}\,Q_{ap^{s}}
  \pmod{p^{3s}},
\]
of depth exactly $3s$ (Finding~\ref{find:weight5}).

\paragraph{(C) Strata and the vanishing defect (\S\ref{sec:strata}).}
With $L=\floor{\log_p n}$ and over \emph{all} $n$, not merely towers,
\[
  \vp(Q_n)\ge0,\qquad
  \vp(\Ph_n/Q_n)=-3L+O(1),\qquad
  \vp(P_n/Q_n)=-5L+O(1),
\]
the $O(1)$ lying in $[-2,1]$ uniformly (Finding~\ref{find:strata}).  Hence the
two Brown--Zudilin linear forms $\Ih_n=Q_n\zeta(3)-\Ph_n$ and
$\Ip_n=Q_n\zeta(5)-P_n$ sit in different valuation strata,
$\vp(\Ih_n/\Ip_n)=2L+O(1)\to+\infty$, and the $p$-adic purity defect vanishes:
\[
  c_p:=\lim_{n\to\infty}\Ih_n/\Ip_n=0
  \qquad\text{in }\Qp,\quad p\in\{5,7,11,13,17,19\}
\]
(Finding~\ref{find:cp}).  Archimedeanly the same limit is
$c=-1/(2\zeta(2))=-3/\pi^2\ne0$, and that non-vanishing is exactly what exiles
the Brown--Zudilin class from the minimal ray \cite{PaperB-inprep}.  The
contrast is structural and is discussed in \S\ref{ss:contrast}: the archimedean
side is obstructed by \emph{positivity}, the $p$-adic side is unobstructed
because the Frobenius \emph{grading} $\mathrm{diag}(1,p^3,p^5)$ separates the
rows.

\paragraph{(D) The exclusion theorem (\S\ref{sec:exclusion}).}
This is a headline negative and we state it as one.  Using an
implementation of the Kubota--Leopoldt $p$-adic $L$-function built from scratch
and validated on $980/980$ interpolation identities (Appendix~\ref{app:KL}):
\begin{itemize}[leftmargin=1.5em]
\item No relation $c_0+c_1\zeta_p(w)+c_2\Lambda=0$ of height $\lesssim10^4$
  holds, for $\Lambda$ any of $\Lambda_a$ (Ap\'ery), $\Lap_a$, $\Lah_a$
  (Brown--Zudilin), at $p=5,7,11,13$, to $12$--$27$ certified $p$-adic digits
  (Finding~\ref{find:exclusion}).  Every candidate returned by exact LLL sits
  \emph{at} the noise floor.
\item The sharp form of the negative is not about margins.  For all $n\le5000$
  and $p\in\{5,7,11,13\}$,
  \begin{gather*}
    \vp\bigl(b_n-\zeta_p(3)a_n\bigr)-\vp(a_n)=-3\floor{\log_pn}+O(1),\\
    \vp\bigl(Q_n\zeta_p(5)-P_n\bigr)-\vp(Q_n)=-5\floor{\log_pn}+O(1)
  \end{gather*}
  (Finding~\ref{find:diverge}).  The $p$-adic linear forms \emph{diverge
  polynomially} where a Bel-type irrationality criterion
  \cite[Lemma~4]{LSZ2025} needs $\exp(-\alpha n)$.  The gap is infinite, not
  finite.
\item Structurally: the Ap\'ery operator
  $\theta^3-t(34\theta^3+51\theta^2+27\theta+5)+t^2(\theta+1)^3$ has order $3$,
  so in the Beukers--Vlasenko Frobenius expansion only $\alpha_1,\alpha_2$
  exist --- and by \cite[Thm.~1.5]{BV2025} both are $0$.  The value
  $\zeta_p(3)$ first occupies $\alpha_3$, which requires order $\ge4$
  (\S\ref{ss:BVposition}).
\end{itemize}

\paragraph{(E) The identification (\S\ref{sec:identify}).}
The exclusions are then \emph{explained}, not merely observed.  Ap\'ery's
formula splits $b_n/a_n$ into a harmonic part and a remainder, and both parts
have exact $p$-adic limits:
\begin{align*}
  h_p(a)&:=\lim_s p^{3s}H_3(ap^s)
        =\sum_{x\in\ZZ[1/p],\ 0<x\le a}x^{-3} &&\text{(Theorem~\ref{thm:hp})},\\
  c_p(A,B)&:=\lim_{t\to\infty}\bin{Ap^t}{Bp^t}
        =-p^{\vp\binom{A}{B}}\,T(A,B)\,
          \frac{\Gp(A+1)}{\Gp(B+1)\Gp(A-B+1)}\,
          e^{\,\mathcal{Z}(A,B)} &&\text{(Theorem~\ref{thm:kaz})},
\end{align*}
%% FIXED (MINOR): T is a product over the base-$p$ TRUNCATIONS $\lfloor A/p^l\rfloor$, not over
%% FIXED the base-$p$ digits.
where $T$ is an explicit finite $\Gp$-product over the base-$p$ truncations
$\floor{A/p^l}$, $\floor{B/p^l}$, $\floor{(A-B)/p^l}$ ($l\ge1$), and
\[
  \mathcal{Z}(A,B)=-\sum_{\substack{m\ge3\\ m\ \mathrm{odd}}}
    \zeta_p(m)\,\bigl(A^m-B^m-(A-B)^m\bigr)\frac{p^m}{m(1-p^m)} .
\]
Assembling, $\Lambda_a=h_p(a)+(\text{convergent block sum})/\Atil_a$ with every
block a polynomial in the $c_p(\cdot,\cdot)$ (Finding~\ref{find:structure}).
Since each $c_p(A,B)$ is the exponential of an infinite $\QQ$-combination of
\emph{all} of $\zeta_p(3),\zeta_p(5),\zeta_p(7),\dots$, no finite low-height
$\QQ$-relation with $\zeta_p(3)$ alone can exist.  The quantitative form of
this: truncating the series after $\zeta_p(3)$ reproduces $4$--$5$ digits out of
$24$--$30$ \VER{$p=5,7,11,13$}.  The same mechanism refutes the hypothesis that
the family $\{\Lambda_a\}$ is generated over $\QQ$ by a single constant
(\S\ref{ss:onegen}).

\paragraph{(F) Existence without archimedean limits (\S\ref{sec:noarch}).}
Three of the fifteen families --- $B$, $(\delta)$, $(\eta)$ --- have
characteristic roots that are complex conjugates of equal modulus, so
$B(n)/A(n)$ oscillates forever and there is \emph{no} archimedean Ap\'ery
limit.  All three nevertheless have $p$-adic tower limits: $(\delta)$ outright
at rate $3$, and $B$, $(\eta)$ after the $\chi(p)^{-s}$ twist.  The $p$-adic
Ap\'ery limit exists where the archimedean one does not.

\subsection{Method, and how to read the evidence tags}\label{ss:evidence}

Statements in this paper carry an explicit evidence class, printed with the
statement.

\begin{description}[leftmargin=2.2em,style=nextline]
\item[\PROVED] A complete proof is given here or in the cited companion paper.
\item[\VER{range}] A finite exact computation over the stated range, at the
  stated $p$-adic precision.  This is evidence, never proof, and the range is
  always stated.
\item[\RECALLED] Asserted from memory of the literature and not checked against
  a source in the preparation of this paper.
\item[\FIT] A numerical fit reproducing known data points, extrapolated.
\end{description}

Statements labelled \textsc{finding} are supported by \VER{} evidence only;
statements labelled \textsc{theorem}, \textsc{proposition}, \textsc{lemma},
\textsc{corollary} are \PROVED.  We have been deliberate about not promoting one
class to the other: several of the most interesting statements below (the rate
law, the strata bounds, $c_p=0$) are at present Findings.

% ---------------------------------------------------------------------
% FLAG FOR RIVER:  this is the ONLY methodology sentence in the paper.
% Keep, rewrite or delete as you prefer; nothing else depends on it.
% ---------------------------------------------------------------------
\begin{remark}[Methodology]\label{rem:method}
All numerical statements below were produced by exact integer and rational
arithmetic --- no floating point occurs anywhere in the pipeline, and every
$p$-adic digit is obtained from an exact rational --- and each closed form
recorded here was checked against an independent direct computation of the
object it describes, written separately from the code that produced the form.
Ranges, certified precisions and the number of independent checks are stated
with each claim; the programs are listed in Appendix~\ref{app:repro}.
\end{remark}

\subsection{Notation and conventions}\label{ss:notation}

$p$ denotes an odd prime, always $\ge5$ unless stated otherwise.  We write
$q_p=p$ (and $q_2=4$), $\omega$ for the Teichm\"uller character modulo $q_p$,
and $\langle x\rangle=x/\omega(x)$.  Following \cite[Def.~8]{LSZ2025} we set
\[
  \zeta_p(s):=L_p\bigl(s,\omega^{1-s}\bigr)\qquad(s\in\ZZ,\ s\ge2),
\]
so $\zeta_p(3)=L_p(3,\omega^{-2})$ and $\zeta_p(5)=L_p(5,\omega^{-4})$; this is
the same normalisation as \cite{BV2025}, for which $\zeta_p(m)$ is the $p$-adic
limit of $-(1-p^{n-1})B_n/n$ along $n=1-m+(p-1)p^r$.  In particular
$\zeta_p(2k)=0$ for all $k\ge1$, a vanishing that will be used repeatedly.

For the classical Ap\'ery pair we write $f(n)=b_n/a_n$, and we compute with the
integer normalisation $\alpha_n=(n!)^3a_n$, $\beta_n=(n!)^3b_n$, which satisfy
\begin{equation}\label{eq:int-rec}
  u_{n+1}=(34n^3+51n^2+27n+5)u_n-n^6u_{n-1},\qquad
  (\alpha_0,\alpha_1)=(1,5),\quad(\beta_0,\beta_1)=(0,6),
\end{equation}
so that $f=\beta/\alpha$.  \VER{$\alpha_n$ against the closed-form double sum for
$n\le29$, $0$ mismatches; $f(400)$ against $\zeta(3)$ to $79$ decimals;
generated exactly to $n=125\,000$, and mod $p^K$ to $n=400\,000$.}

For the Brown--Zudilin pair we use the printed normalisation of
\cite{BrownZudilin2022}: $Q_0=1$, $P_0=0$, $Q_1=21$, $P_1=87/4$,
$P_2=1190161/384$, $\Ph_1=101/4$, $\Ph_2=344923/96$, and we set
$\Ip_n=Q_n\zeta(5)-P_n$, $\Ih_n=Q_n\zeta(3)-\Ph_n$.  \VER{the certified order-$3$
recurrence has residual exactly $0$ on all three rows for every $n\le357$
($358$ checks per row), and the exact extension reproduces
$P_n/Q_n\to\zeta(5)$, $\Ph_n/Q_n\to\zeta(3)$ to $59$ decimals at $n=800$ and
$n=5000$.}

Throughout $L=\floor{\log_p n}$ and $q=\floor{n/p}$.

\subsection{Organisation}
Section~\ref{sec:engine} recalls the descent congruences and turns them into a
convergence engine.  Section~\ref{sec:rate} is the rate law and its weight-$5$
correction.  Section~\ref{sec:strata} is the strata theorem and $c_p=0$.
Section~\ref{sec:exclusion} is the exclusion theorem.
Section~\ref{sec:identify} identifies the limits.  Section~\ref{sec:noarch} is
the no-archimedean-limit trio.  Section~\ref{sec:open} lists open problems.
Appendix~\ref{app:KL} documents and validates the Kubota--Leopoldt
implementation; Appendix~\ref{app:LSZ} reproduces the Lai--Sprang--Zudilin
denominator ledger; Appendix~\ref{app:data} collects numerical data;
Appendix~\ref{app:repro} lists the programs.

% =====================================================================
\section{The descent congruences as the convergence engine}\label{sec:engine}
% =====================================================================

\subsection{The sporadic families}\label{ss:families}

Table~\ref{tab:families} lists the fifteen sporadic Ap\'ery-like pairs, in the
labelling of \cite{MalikStraub}.  The recurrences are
\begin{align}
  (\mathrm{R}2)&\quad (n+1)^2u_{n+1}=(an^2+an+b)u_n-cn^2u_{n-1},\label{eq:R2}\\
  (\mathrm{R}3)&\quad (n+1)^3u_{n+1}=(2n+1)(an^2+an+b)u_n-n(cn^2+d)u_{n-1},\label{eq:R3}
\end{align}
with $A$ the sporadic binomial sum and $B$ the second solution normalised by
$B(0)=0$, $B(1)=1$ (the recurrence being enforced for $n\ge1$; the $n=0$ step is
degenerate --- this is Ap\'ery's own construction, and the normalisation of
\cite{ChamberlandStraub}).  Rescaling $B$ by a $p$-unit
rational changes $\Lambda_a$ only by that scalar, so for $p\ge5$ this
normalisation is canonical up to a harmless constant.

\begin{table}[t]
\centering\small
\caption{The fifteen sporadic pairs, their Ap\'ery limits, weights and
characters, after \cite{MalikStraub,Straub2023}; limits identified in
\cite{PaperA-inprep}, several of them presumably also in the numerical tables of
\cite{ASvSZ2008}.  $w$ is the minimal exponent with $d_n^wB(n)\in\ZZ$.}
\label{tab:families}
\begin{tabular}{@{}llllcc@{}}
\toprule
\# & label & $A(n)$ & Ap\'ery limit & $w$ & $\chi$\\
\midrule
1 & $A$ & $\sum_k\bin{n}{k}^3$ (Franel) & $\zeta(2)/4$ & 2 & $1$\\
2 & $B$ & $\sum_k(-1)^k3^{n-3k}\bin{n}{3k}\frac{(3k)!}{k!^3}$ & \emph{none} & 2 & $\chi_{-3}$\\
3 & $C$ & $\sum_k\bin{n}{k}^2\bin{2k}{k}$ & $L_{-3}(2)/2$ & 2 & $\chi_{-3}$\\
4 & $D$ & $\sum_k\bin{n}{k}^2\bin{n+k}{n}$ (Ap\'ery $\zeta(2)$) & $\zeta(2)/5$ & 2 & $1$\\
5 & $E$ & $\sum_k\bin{n}{k}\bin{2k}{k}\bin{2(n-k)}{n-k}$ & $L_{-4}(2)/2=G/2$ & 2 & $\chi_{-4}$\\
6 & $F$ & $\sum_k(-1)^k8^{n-k}\bin{n}{k}\sum_l\bin{k}{l}^3$ & $(5/8)L_{-3}(2)$ & 2 & $\chi_{-3}$\\
7 & $(\alpha)$ & $\sum_k\bin{n}{k}^2\bin{2k}{k}\bin{2(n-k)}{n-k}$ (Domb) & $7\zeta(3)/24$ & 3 & $1$\\
8 & $(\gamma)$ & $\sum_k\bin{n}{k}^2\bin{n+k}{n}^2$ (Ap\'ery $\zeta(3)$) & $\zeta(3)/6$ & 3 & $1$\\
9 & $(\delta)$ & $\sum_k(-1)^k3^{n-3k}\bin{n}{3k}\bin{n+k}{n}\frac{(3k)!}{k!^3}$ & \emph{none} & 3 & $1$\\
10 & $(\varepsilon)$ & $\sum_k\bin{n}{k}^2\bin{2k}{n}^2$ & $7\zeta(3)/32$ & 3 & $1$\\
11 & $(\zeta)$ & $\sum_{k,l}\bin{n}{k}^2\bin{n}{l}\bin{k}{l}\bin{k+l}{n}$ & $L_{-3}(3)/3$ & 3 & $\chi_{-3}$\\
12 & $(\eta)$ & $\sum_k(-1)^k\bin{n}{k}^3\bigl[\bin{4n-5k-1}{3n}+\bin{4n-5k}{3n}\bigr]$ & \emph{none} & 3 & $\chi_{5}$\\
13 & $s_7$ & $\sum_k\bin{n}{k}^2\bin{n+k}{k}\bin{2k}{n}$ & $\zeta(2)/7$ & 2 & $1$\\
14 & $s_{10}$ & $\sum_k\bin{n}{k}^4$ & $\zeta(2)/5$ & 2 & $1$\\
15 & $s_{18}$ & Cooper's $s_{18}$ & $L_{-3}(2)/2$ & 2 & $\chi_{-3}$\\
\bottomrule
\end{tabular}
\end{table}

\subsection{The twisted descent law}\label{ss:LB}

The input we take from Paper~A \cite{PaperA-inprep} is the following universal
digit-descent law.  Write $q=\floor{n/p}$.

\begin{theorem}[{Paper A \cite{PaperA-inprep}}]\label{thm:LB}
\PROVED\ (single-digit form).  Let $p\ge5$, let $(A,B)$ be a sporadic pair with
weight $w$ and character $\chi$ as in Table~\ref{tab:families}, and suppose $B$
admits a $\chi$-homogeneous harmonic-monomial decomposition of weight $w$ on the
summand of $A$ satisfying the tameness and Lucas hypotheses of
\cite{PaperA-inprep}.  Then for $1\le a<p$, $0\le r<p$ and $n=ap+r$,
\[
  p^{w}\,B(ap+r)\ \equiv\ \chi(p)^{e}\,B(a)\,A(r)\pmod p ,
\]
with $e\in\{0,1\}$ the number of $\chi$-letters in each monomial.  The
hypotheses are verified there for $(\gamma)$, $A$, $s_{10}$ $(p\ge5$, resp.\
$p\ge7)$ and $D$ $(p\ge7)$.
\end{theorem}

\begin{finding}[{the master form, Paper~A}]\label{find:master}
\VER{all $15$ sporadic pairs; $28$ primes $5\le p\le103$; $n\le400$ for all $p$
and $n\le1200$ for $p=5,7,11$; exact rational arithmetic; zero failures}
\[
  (\mathrm{LB}_w^\chi)\qquad
  \vp\bigl(p^{w}B_nA_q-\chi(p)B_qA_n\bigr)\ \ge\ w,
  \qquad q=\floor{n/p},\qquad p\ge5,
\]
and the floor is \emph{exactly} $w$ in every single (sequence, prime) cell: no
$\kappa$-dependence, no centre-residue exception, no exceptional-digit
exception.  For the classical Ap\'ery pair this reads $p^3b_na_q\equiv b_qa_n
\pmod{p^3}$, \VER{$5\le p\le31$, $n\le320$}.
\end{finding}

The mechanism, which we will need again in \S\ref{ss:harmrate}, is a two-line
lemma about harmonic letters.

\begin{lemma}[{Paper A \cite{PaperA-inprep}}]\label{lem:K}
\PROVED.  Let $\chi$ be a completely multiplicative arithmetic function
($\chi\equiv1$ allowed), $r\ge1$, $y\ge0$, and
$K^{(r)}_\chi(y)=\sum_{m\le y}\chi(m)m^{-r}$.  Then, as an identity in $\QQ$,
\[
  p^{r}K^{(r)}_\chi(y)\ =\ \chi(p)\,K^{(r)}_\chi(\floor{y/p})\ +\ p^{r}T,
  \qquad T=\!\!\sum_{m\le y,\ p\nmid m}\!\!\chi(m)m^{-r}\in\ZZ_{(p)} .
\]
\end{lemma}

\begin{proof}
Split the sum at $p\mid m$.  The $p$-divisible part is
\[
  \sum_{j\le\floor{y/p}}\chi(jp)(jp)^{-r}
  =\chi(p)\,p^{-r}\!\!\sum_{j\le\floor{y/p}}\!\!\chi(j)j^{-r},
\]
and the remaining terms are $p$-integral.
\end{proof}

The factor $\chi(p)$ in $(\mathrm{LB}_w^\chi)$ is exactly the $\chi(p)$ produced
by the $p$-divisible layer of the harmonic weight.  It is invisible in the
$A$-row: the unit-root congruence $A(ap+r)\equiv A(a)A(r)\pmod p$ is untwisted
for all fifteen families \cite{Gessel1982,MalikStraub}.  The character
appears only on the non-unit Frobenius eigenvalue.

\subsection{Existence of the tower limits}\label{ss:exists}

\begin{theorem}\label{thm:exists}
\PROVED, given $(\mathrm{LB}_w^\chi)$ in ratio form at every level of the tower.
Let $p\ge5$, $p\nmid a$, $n_s=ap^s$, and suppose $A(n_s)$ is a $p$-unit cell for
all $s$, so that the ratio form
\[
  p^{w}f(n_s)\ \equiv\ \chi(p)\,f(n_{s-1})\pmod{p^{w}},
  \qquad f=B/A,
\]
holds.  Put $L_s:=\chi(p)^{s}p^{ws}f(n_s)$.  Then
\[
  L_s\equiv L_{s-1}\pmod{p^{ws}},
\]
so $(L_s)_{s\ge0}$ is Cauchy in $\Qp$ and
$\Lambda_a:=\lim_{s\to\infty}L_s$ exists, with
$\Lambda_a\equiv L_s\pmod{p^{w(s+1)}}$.  The same holds along any descent branch
$n_s=pn_{s-1}+r_s$.
\end{theorem}

\begin{proof}
$L_s-L_{s-1}=\chi(p)^{s}p^{w(s-1)}\bigl(p^{w}f(n_s)-\chi(p)f(n_{s-1})\bigr)$,
and the bracket has valuation $\ge w$; hence $\vp(L_s-L_{s-1})\ge ws$.  The
sequence is therefore Cauchy, and telescoping the tail gives
$\vp(\Lambda_a-L_s)\ge w(s+1)$.  Nothing in the argument used $n_s=ap^s$ beyond
$n_{s-1}=\floor{n_s/p}$.
\end{proof}

For the classical Ap\'ery pair ($w=3$, $\chi=1$) this is
\[
  L_s=p^{3s}\frac{b_{ap^s}}{a_{ap^s}},\qquad
  L_s\equiv L_{s-1}\pmod{p^{3s}},\qquad \Lambda_a=\lim_sL_s .
\]
For the Brown--Zudilin rows we write, in the same spirit,
\[
  \Lap_a:=\lim_s p^{5s}\frac{P_{ap^s}}{Q_{ap^s}},
  \qquad
  \Lah_a:=\lim_s p^{3s}\frac{\Ph_{ap^s}}{Q_{ap^s}} .
\]
For the $P$-row Theorem~\ref{thm:exists} does \emph{not} apply as stated,
because the naive weight-$5$ ratio descent degrades; see \S\ref{ss:weight5}.

\subsection{What the limit does and does not remember}\label{ss:invariance}

\begin{proposition}\label{prop:invariance}
\PROVED.  In the situation of Theorem~\ref{thm:exists}:
\begin{enumerate}[label=(\roman*),leftmargin=2em]
\item $\Lambda_a$ is unchanged by $B\mapsto B+cA$ for any $c\in\QQ$;
\item $\Lambda_a$ scales by $\lambda$ under $B\mapsto\lambda B$;
\item $\Lambda_{p^ka}=\chi(p)^{k}p^{-wk}\Lambda_a$ for $k\ge0$ (so
  $\Lambda_{p^ka}=p^{-wk}\Lambda_a$ whenever $\chi$ is trivial);
%% FIXED (CRITICAL): item (iv) was FALSE over the range its own \VER tag claimed. The
%% FIXED strengthening "in fact modulo p^{w+v_p(Lambda_a)}" has been DELETED (it fails at 4 of
%% FIXED the 36 p-unit cells, and 8 of 44 cells overall), and the p-unit-cell hypothesis --- in
%% FIXED force via Theorem \ref{thm:exists} but silently overridden by the old \VER range --- is
%% FIXED now stated inside the item. Independently recomputed here (exact Fraction Apery
%% FIXED recurrence, s=0,1,2): failures of the mod-p^w half are exactly the 8 cells with
%% FIXED v_p(A(a))>0, i.e. exactly the non-p-unit cells. The old claim was inherited from
%% FIXED LAMBDA_HUNT:114 / PADIC_SEAM:122, which are themselves wrong.
\item $\Lambda_a\equiv f(a)\pmod{p^{\,w}}$, \emph{provided} $A(a)$ is a $p$-unit cell (which is
  a standing hypothesis of Theorem~\ref{thm:exists}, and is not automatic: see
  Remark~\ref{rem:ivsharp}).
\end{enumerate}
\end{proposition}

\begin{proof}
(i) The perturbation contributes $\chi(p)^sp^{ws}c\to0$.  (ii) is immediate.
(iii) Reindexing $t=s+k$ in
$\Lambda_{p^ka}=\lim_s\chi(p)^{s}p^{ws}f(ap^{s+k})$ gives
$\chi(p)^{-k}p^{-wk}\Lambda_a$, and $\chi(p)^{-k}=\chi(p)^{k}$ for a quadratic
(or trivial) character at an unramified prime.  (iv) is the case $s=0$ of
$\Lambda_a\equiv L_s\pmod{p^{w(s+1)}}$ in Theorem~\ref{thm:exists}.
\end{proof}

Item~(i) is the reason $\Lambda_a$ is an intrinsic invariant of the recurrence
and the tower rather than of the particular second solution --- and equally the
reason it cannot ``see'' the archimedean normalisation $\lim_nb_n/a_n=\zeta(3)$
except through the single global scalar.  This is the first structural sign that
the expectation $\Lambda\leftrightarrow\zeta_p(3)/6$ is misplaced; \S\ref{sec:exclusion}
makes it a theorem-strength negative and \S\ref{sec:identify} explains it.  Item~(iii)
is used throughout as an internal consistency check on the numerics.  Item~(iv)
is the ``trivial floor'': on a $p$-unit cell, any ansatz expressing $\Lambda_a$ through $f(a)$
reproduces $w$ digits for free, so only agreement substantially beyond $w$
digits is evidence of anything.  This is the calibration used in
\S\ref{ss:onegen}.

\begin{remark}[how sharp (iv) is, and where it stops]\label{rem:ivsharp}
Both qualifications matter, and we record the sweep because the floor is a calibration the rest
of the paper leans on.  For the Ap\'ery pair ($w=3$) over the $44$ cells $p\in\{5,7,11,13\}$,
$a\le12$, $p\nmid a$:
\begin{itemize}[leftmargin=1.4em]
\item $36$ cells are $p$-unit cells, and (iv) holds at every one of them.  It is
  \emph{sharp}: $\vp(\Lambda_a-f(a))=w=3$ exactly at $26$ of the $36$, with $4$ at $\vp=4$,
  $2$ at $5$ and $1$ at $6$ --- so nothing beyond $w$ digits is free.
\item The remaining $8$ cells --- $p=5$, $a=1,3,6,7,8,9,11$ and $p=11$, $a=5$ --- have
  $\vp(A(a))>0$, and (iv) \emph{fails} at every one of them: the difference-valuation is $1$ at
  $p=5$, $a=1,3,7,9$ and at $p=11$, $a=5$; $-1$ at $p=5$, $a=6,8$; and $-4$ at $p=5$, $a=11$.
  These are exactly the cells excluded by Theorem~\ref{thm:exists}'s hypothesis.
\item The natural strengthening to $p^{\,w+\vp(\Lambda_a)}$ is \emph{false} even on $p$-unit
  cells: it fails at $p=7$, $a=3$ and at $p=13$, $a=2,6,10$, where $\vp(\Lambda_a)=1$ so the
  claim would demand $p^{4}$ while the true difference-valuation is exactly $3$.
\end{itemize}
The simplest counterexample needs no computation beyond Table~\ref{tab:constants}: at $p=5$ it
prints $\Lambda_1=5^{-1}(1\,1\,3\,3\,4\dots)$, while $f(1)=b_1/a_1=6/5=5^{-1}(1\,1\,0\,0\,0\dots)$;
they first differ in the third digit, so $v_5(\Lambda_1-f(1))=1$, which is neither $\ge w=3$ nor
$\ge w+v_5(\Lambda_1)=2$.  (Here $a_1=5$, so $a=1$ is not a $5$-unit cell.)
\end{remark}

% =====================================================================
\section{The rate law}\label{sec:rate}
% =====================================================================

Theorem~\ref{thm:exists} gives a rate of at least $w$ new $p$-adic digits per
tower level.  The measured rate is neither $w$ nor constant; it is
$\min(3,r(\chi,w))$, and both arguments of the minimum have a mechanism.

\subsection{The harmonic tower and its rate}\label{ss:harmrate}

Let $\chi$ be a Dirichlet character (possibly trivial), $w\ge2$, and set
$K^{(w)}_\chi(y)=\sum_{m\le y}\chi(m)m^{-w}$.  By Lemma~\ref{lem:K} the tower
\[
  h^{\chi,w}_p(a)\ :=\ \lim_{s\to\infty}\chi(p)^{s}p^{ws}K^{(w)}_\chi(ap^s)
\]
exists, and its convergence rate is computable from the layer sums.

\begin{finding}\label{find:harmrate}
\VER{$p=5,7,11$; $w=2,3$; $a\le3$; $s\le3$} The harmonic tower converges at the
rates of Table~\ref{tab:harmonicrate}.
\end{finding}

\begin{table}[h]
\centering\small
\caption{Rate of the $\chi$-twisted harmonic tower, and the mechanism.}
\label{tab:harmonicrate}
\begin{tabular}{@{}ccll@{}}
\toprule
$\chi$ & $w$ & rate $r(\chi,w)$ & mechanism\\
\midrule
$1$ & $2$ & $3$ & $S_i\approx B_Mp^{i}$; $\zeta_p(3)\ne0$\\
$1$ & $3$ & $5$ (digits $4,9,14,19$ at $p=5$) & idem, but $\zeta_p(4)=0$ kills the leading term\\
$\chi\ne1$ & $2$ & $2$ & $\sum\chi(k)k^{-w}$ tends to a \emph{nonzero} $L$-value: no $p^{i}$ gain\\
$\chi\ne1$ & $3$ & $3$ & idem\\
\bottomrule
\end{tabular}
\end{table}

In words: for the trivial character the leading layer of the tower is divisible
by an extra power of $p$, gaining one digit; for odd $w$ it gains two, because
the coefficient that would have carried the first correction is
$\zeta_p(w+1)=0$.  For a nontrivial character the layer sums converge to a
nonzero generalised-Bernoulli/$L$-value and there is no gain at all.  Thus
\[
  r(\chi,w)=w\ (\chi\ne1),\qquad
  r(1,w)=\begin{cases} w+1,& w\ \text{even},\\ w+2,& w\ \text{odd}.\end{cases}
\]
\VER{$w=2,3$ only.}

\subsection{The Jacobsthal--Kazandzidis cap}\label{ss:cap}

The second argument of the minimum is a cap, and it comes from binomial
coefficients rather than from harmonic sums.

\begin{finding}\label{find:kazrate}
\VER{$p=5,7,11,13$; $12$ pairs $(A,B)$ per prime; $t\le5$} The
Jacobsthal--Kazandzidis tower $\bin{Ap^t}{Bp^t}$ converges at exactly $3$
digits per level \cite{JacobsthalKazandzidis}.
\end{finding}

This is weight-independent.  Since (\S\ref{sec:identify}) the non-harmonic part
of $\Lambda_a$ is assembled entirely out of the limits $c_p(A,B)$ of such
towers, no tower limit built this way can converge faster than $3$ digits per
level, whatever its weight.

\subsection{The rate law}\label{ss:ratelaw}

\begin{finding}[Rate law]\label{find:ratelaw}
\VER{$15$ sporadic families $\times$ $p\in\{5,7,11\}$, $N=20\,000$, towers
$a=1,2$} The number of new certified $p$-adic digits gained per level of the
tower $n_s=ap^s$ is
\[
  \min\bigl(3,\ r(\chi,w)\bigr),
\]
%% FIXED (MAJOR M-4): the exception list omitted F, which at p=5 deviates in exactly the same
%% FIXED way and by exactly the same amount as C and s_18 (independent reimplementation from the
%% FIXED integer recurrence for (n!)^k u_n: at p=5, cumulative agreement C = F = s18 =
%% FIXED 4,5,7,9,11,13 against a predicted 2,4,6,8,10,12, while B and E are on the nose).
%% FIXED Inherited from LAMBDA_HUNT:291, not a transcription slip. Also fixed: the deviations are
%% FIXED upward only cumulatively -- the p=5 level-2 increment for C, F, s18 is 1, below the
%% FIXED predicted 2, and the p=7 increments for (epsilon) are 4, 5, 2, 3, 3.
with $r(\chi,w)$ as in Table~\ref{tab:harmonicrate}, in every cell except at
$(\varepsilon)$, $(\zeta)$, $C$, $F$, $s_{18}$, where
blocks cancel.  In particular the rate is $3$ for all eight trivial-character
families at both weights, and $2$ for the $\chi$-twisted weight-$2$ families
$B$, $C$, $E$, $F$, $s_{18}$.
\end{finding}

\begin{remark}[the shape of the deviations]
They are \emph{cumulative} excesses, not uniformly faster levels.  At $p=5$ the three families
$C$, $F$ and $s_{18}$ behave identically --- cumulative agreement $4,5,7,9,11,13$ against the
predicted $2,4,6,8,10,12$ --- so the excess is $+2$ acquired at the first level and then held,
and the level-$2$ increment is $1$, i.e.\ \emph{below} the predicted $2$.  ($B$ and $E$ are on
the nose at $p=5$, and $F$ is on the nose at $p=7$, where instead $C$ and $s_{18}$ carry a $+1$.)
At $p=7$ the $(\varepsilon)$ increments are $4,5,2,3,3$, one of which is again below the
predicted $3$.  Everything not named in Finding~\ref{find:ratelaw} reproduces
$\min(3,r(\chi,w))$ exactly.
\end{remark}

Table~\ref{tab:familyrates} records the measurements.

\begin{table}[h]
\centering\small
\caption{Digits gained per tower level, $a=1$.  ``raw'' is without the
$\chi(p)^{s}$ twist; ``tw.'' is with it.  $\chi(p)=-1$ cells are the ones where
the raw tower does not converge at all.}
\label{tab:familyrates}
\begin{tabular}{@{}llccccccc@{}}
\toprule
family & $w$ & $\chi$ & $\chi(7)$ & raw ($p=7$) & tw.\ ($p=7$) & tw.\ ($p=11$) & tw.\ ($p=5$)\\
\midrule
$A$ (Franel) & 2 & $1$ & $+1$ & 3 & 3 & 3 & 3\\
$D$ (Ap\'ery $\zeta(2)$) & 2 & $1$ & $+1$ & 3 & 3 & 3 & 3\\
$s_7$, $s_{10}$ & 2 & $1$ & $+1$ & 3 & 3 & 3 & 3\\
$(\alpha)$, $(\gamma)$, $(\delta)$, $(\varepsilon)$ & 3 & $1$ & $+1$ & 3 & 3 & 3 & 3\\
$B$, $C$, $F$ & 2 & $\chi_{-3}$ & $+1$ & \textbf{2} & 2 & 2 & 2 (twist needed)\\
$s_{18}$ & 2 & $\chi_{-3}$ & $+1$ & 2 & 2 & 2 & 2 (twist needed)\\
$E$ & 2 & $\chi_{-4}$ & $\mathbf{-1}$ & \textbf{0} & \textbf{2} & 2 & 2 ($\chi(5)=+1$)\\
$(\zeta)$ & 3 & $\chi_{-3}$ & $+1$ & 4 & 4 & 4 & 4 (twist needed)\\
$(\eta)$ & 3 & $\chi_{5}$ & $\mathbf{-1}$ & \textbf{0} & \textbf{3} & 3 & degenerate\\
\bottomrule
\end{tabular}
\end{table}

\begin{remark}\label{rem:not-weight}
A first pass over weight-$3$ and weight-$5$ data suggests the clean statement
``the rate is always exactly $3$, independently of the weight''.  That statement
is an artefact of testing only trivial-character pairs, where
$\min(3,r(1,w))=3$ for $w\ge2$.  Finding~\ref{find:ratelaw} is the corrected
form.  The moral is worth stating plainly: \emph{the $3$ is
Jacobsthal--Kazandzidis, not the motivic weight}.
\end{remark}

For the classical Ap\'ery pair the rate is flat to a degree worth recording
separately.

\begin{finding}\label{find:aperyrate}
%% FIXED (MAJOR M-3): two universal quantifiers were added on transcription and are false.
%% FIXED (1) "constant in s for every tower and equals v_p(b_a/a_a)" fails at p=5, a=11:
%% FIXED recomputed exactly here, v_p(p^{3s}f(11*5^s)) = -4, -2, -2, -2 for s=0..3 while
%% FIXED v_p(b_11/a_11) = -4. The source logs already showed it (hunt_p5.txt "a=11 v(L)=-2";
%% FIXED t1_p5.log "a=11 ... relagree L:[0, 2, 5, 8, ...]"). PADIC_SEAM T1.2, whence the
%% FIXED sentence, tested only a=1,2 (p=5), a=1,3 (p=7), a=1 (p=11,13).
%% FIXED (2) "flat in the descent branch" is contradicted by Table \ref{tab:aperyrate} itself:
%% FIXED the row "7, branch r=1" reads 4, 6, 9, 13, 15, increments 4, 2, 3, 4, 2.
\VER{exact data to $n=125\,000$: $p=5$, $s\le7$; $p=7$, $s\le6$; $p=11,13$,
$s\le4$} For the Ap\'ery $\zeta(3)$ pair the rate is $3$ digits per
level asymptotically, flat in $p$ and in the tower base $a$; along a $p$-power tower
$n_s=ap^s$ the valuation $\vp\bigl(p^{3s}f(n_s)\bigr)$ is eventually constant in $s$, so that
the pole order of $b_n/a_n$ along a tower is exactly $3\floor{\log_pn}+O(1)$.  See
Table~\ref{tab:aperyrate}.
\end{finding}

\begin{remark}[two qualifications that the data forces]\label{rem:notflat}
Neither ``constant in $s$'' nor ``flat in the branch'' can be strengthened to a universal
statement.
\begin{itemize}[leftmargin=1.4em]
\item \emph{The valuation can jump once, at the first level.}  At $p=5$, $a=11$ we compute
  $\vp\bigl(5^{3s}f(11\cdot5^s)\bigr)=-4,-2,-2,-2$ for $s=0,1,2,3$, while
  $v_5(b_{11}/a_{11})=-4$.  So the constant is $\vp(\Lambda_a)$, which need \emph{not} equal
  $\vp(b_a/a_a)$: here $a_{11}$ is divisible by $5^4$, so $a=11$ is not a $5$-unit cell, and
  Theorem~\ref{thm:exists} does not apply at $s=0$.  The jump is visible in the first entry of
  the corresponding agreement row, which is $0$ rather than $3$.
\item \emph{Along a general descent branch the rate is $3$ only on average.}  The last row of
  Table~\ref{tab:aperyrate} ($p=7$, branch $r=1$) reads $4,6,9,13,15$, i.e.\ per-level
  increments $4,2,3,4,2$.  It is the $p$-power towers, not all branches, on which the increment
  is $3$ level by level.
\end{itemize}
\end{remark}

\begin{table}[h]
\centering\small
\caption{Digits of agreement between $L_s$ and $L_{s-1}$ for the Ap\'ery pair,
after clearing the common $p^{\min v}$ shift.}
\label{tab:aperyrate}
\begin{tabular}{@{}llcl@{}}
\toprule
$p$ & tower & $\vp(\Lambda_a)$ & agreement per level\\
\midrule
5 & $a=1$ & $-1$ & 2, 5, 8, 11, 14, 17, 20\\
5 & $a=2$ & $0$ & 3, 6, 9, 12, 15, 18\\
7 & $a=1$ & $0$ & 3, 6, 9, 12, 15, 18\\
7 & $a=3$ & $1$ & 2, 5, 8, 11, 14\\
11 & $a=1$ & $0$ & 3, 6, 9, 12\\
13 & $a=1$ & $0$ & 3, 6, 9, 12\\
5 & branch $n_s=p^{s+1}-1$ & $1$ & 5, 8, 11, 14, 17, 20\\
7 & branch $r=1$ & $0$ & 4, 6, 9, 13, 15\\
\bottomrule
\end{tabular}
\end{table}

\begin{remark}[an honest precision budget]\label{rem:precision}
Three digits per level is a hard constraint on what can be computed.  Sixty
digits would require $s=19$, i.e.\ $n=p^{19}\approx10^{13}$ at $p=5$; even the
$O(n)$ closed-form-sum route, which needs only $K+O(\log n)$ digits of working
precision rather than the $O(N^2)$ of the recurrence route, caps out near
$n\approx10^7$, i.e.\ about $33$ digits.  The certified precisions actually
reached are in Table~\ref{tab:precision}.  What made the identification of
\S\ref{sec:identify} possible was therefore not more digits but exact closed
forms, which can be evaluated to any precision at any $p$.
\end{remark}

\begin{table}[h]
\centering\small
\caption{Certified relative digits of $\Lambda_a$ (Ap\'ery pair).}
\label{tab:precision}
\begin{tabular}{@{}lcccc@{}}
\toprule
$p$ & $N$ reached & $a=1$ & $a\le4$ & $a\le12$\\
\midrule
5 & $400\,000$ & $s=8$, \textbf{27} & 24 & 21 ($a=12$: $s=6$)\\
7 & $250\,000$ & $s=6$, \textbf{21} & 21--18 & 18\\
11 & $170\,000$ & $s=5$, \textbf{18} & 15 & 12\\
13 & $380\,000$ & $s=5$, \textbf{18} & 15 & 15\\
\bottomrule
\end{tabular}
\end{table}

\subsection{The weight-$5$ correction}\label{ss:weight5}

The Brown--Zudilin $P$-row is the only weight-$5$ instance available to us, and
it is the one where the naive statement fails.

\begin{finding}\label{find:weight5}
\VER{exact ladders to $n=5000$; $p\in\{5,7,11,13\}$}
\begin{enumerate}[label=(\roman*),leftmargin=2em]
\item $\vp\bigl(p^{5s}P_{ap^s}/Q_{ap^s}\bigr)$ is constant in $s$: weight $5$ is
  the right and only normalising exponent for the $P$-row.  Likewise
  $\vp\bigl(p^{3s}\Ph_{ap^s}/Q_{ap^s}\bigr)$ is constant in $s$.
\item The rate is nevertheless $3$, not $5$, for both rows
  (Table~\ref{tab:bzrate}).
\item The single-step ratio descent $p^{5}f(n_s)\equiv f(n_{s-1})$ has depth
  $5-2s$, which \emph{degrades} and goes negative at $s\ge3$.  The multi-digit
  statement that survives is the $p^{5L}$-scaled one,
  \begin{equation}\label{eq:scaled5}
    p^{5s}\,P_{ap^{s}}\,Q_{ap^{s-1}}
    \ \equiv\ p^{5(s-1)}\,P_{ap^{s-1}}\,Q_{ap^{s}}
    \pmod{p^{3s}},
  \end{equation}
  with depth exactly $3s$.
\item The $\Ph$-row, by contrast, has \emph{constant} single-step depth $3$,
  exactly like the Ap\'ery $b$-row.
\end{enumerate}
\end{finding}

\begin{table}[h]
\centering\small
\caption{Brown--Zudilin towers, $a=1$: valuation of the normalised row and
digits of agreement per level.}
\label{tab:bzrate}
\begin{tabular}{@{}llccl@{}}
\toprule
$p$ & row (weight) & $s$ range & $\vp(p^{ws}\cdot\text{row}/Q)$ & digits of agreement\\
\midrule
5 & $P$ (5) & $0..5$ & $0,0,0,0,0,0$ & 3, 6, 9, 12, 15\\
5 & $\Ph$ (3) & $0..5$ & $0,0,0,0,0,0$ & 3, 6, 9, 12, 15\\
7 & $P$ (5) & $0..4$ & $-1$ (const) & 2, 5, 8, 11\\
7 & $\Ph$ (3) & $0..4$ & $-1$ (const) & 2, 5, 8, 11\\
11 & $P$ (5) & $0..3$ & $0$ (const) & 4, 7, 10\\
13 & $P$ (5) & $0..3$ & $0$ (const) & 3, 6, 9\\
\bottomrule
\end{tabular}
\end{table}

Equation~\eqref{eq:scaled5} is the correct form of the ``graded-weight''
descent for the top row.  It is worth emphasising what goes wrong with the
naive form: it is not that the congruence is false at one level, but that its
depth is $5-2s$, so that it is \emph{true and useless} for $s\ge3$.  Only the
$Q$-weighted, $p^{5L}$-scaled statement has a depth that grows with $s$, and
its growth rate is $3s$ --- once again the Kazandzidis $3$, not the weight $5$.

% =====================================================================
\section{Valuation strata and the vanishing of the $p$-adic defect}\label{sec:strata}
% =====================================================================

\subsection{The strata}\label{ss:strata}

\begin{finding}[Strata]\label{find:strata}
\VER{$n\le5000$; $p\in\{5,7,11,13,17,19\}$}
Over \emph{all} $n$, not merely along towers, and with $L=\floor{\log_pn}$,
\[
  \vp(Q_n)\ \ge\ 0,\qquad
  \vp(\Ph_n/Q_n)=-3L+O(1),\qquad
  \vp(P_n/Q_n)=-5L+O(1),
\]
with the $O(1)$ bounded by $2$ uniformly; explicitly the deviations lie in the
sets of Table~\ref{tab:strata}, all contained in $[-2,1]$.
\end{finding}

\begin{table}[h]
\centering\small
\caption{Deviations of the row valuations from the predicted strata, all $n\le5000$.}
\label{tab:strata}
\begin{tabular}{@{}cccc@{}}
\toprule
$p$ & $\vp(\Ph_n/Q_n)+3L$ & $\vp(P_n/Q_n)+5L$ & $\vp(\Ph_n/P_n)-2L$\\
\midrule
5 & $\{-1,0\}$ & $\{0,1\}$ & $\{-1,0\}$\\
7 & $\{-2,-1,0,1\}$ & $\{-2,-1,0,1\}$ & $\{-2,-1,0\}$\\
11 & $\{-1,0\}$ & $\{-1,0,1\}$ & $\{-1,0\}$\\
13 & $\{-1,0,1\}$ & $\{0,1\}$ & $\{-2,-1,0,1\}$\\
17 & $\{-1,0\}$ & $\{0,1\}$ & $\{-1,0\}$\\
19 & $\{-1,0\}$ & $\{-1,0,1\}$ & $\{-2,-1,0,1\}$\\
\bottomrule
\end{tabular}
\end{table}

\begin{remark}\label{rem:Qvals}
$\vp(Q_n)=0$ for every $n\le5000$ at $p=5,13,17$, but reaches $7$ at $p=7$,
$4$ at $p=11$ and $3$ at $p=19$: the exceptional-prime phenomenon is present in
the $Q$-row and must be carried through any argument that divides by $Q_n$.
\end{remark}

Finding~\ref{find:strata} is a clean denominator statement about the whole
Brown--Zudilin family, and it is precisely the input that a Rhin--Viola-style
denominator saving \cite{RhinViola} would have to improve upon.

\subsection{$c_p=0$}\label{ss:cp}

\begin{finding}[the $p$-adic defect vanishes]\label{find:cp}
\VER{$n\le5000$; $p\in\{5,7,11,13,17,19\}$} Since
$\zeta_p(3),\zeta_p(5)\in p^{-1}\Zp$ while the rows have poles of order $3L$
and $5L$,
\[
  \vp(\Ih_n)-\vp(Q_n)=-3L+O(1),\qquad
  \vp(\Ip_n)-\vp(Q_n)=-5L+O(1),
\]
hence $\vp(\Ih_n/\Ip_n)=2\floor{\log_pn}+O(1)\to+\infty$ and
\[
  \boxed{\ c_p\ :=\ \lim_{n\to\infty}\frac{\Ih_n}{\Ip_n}\ =\ 0\quad\text{in }\Qp . \ }
\]
\end{finding}

\begin{finding}[the renormalised defect is not a constant]\label{find:chat}
\VER{$p=5,7,11,13$; $a=1,2$; $12$--$18$ certified digits} Dividing out the
$p^{2s}$,
\[
  \widehat{c}_p(a):=\frac{\Lah_a}{\Lap_a}
  =\lim_{s}p^{-2s}\frac{\Ph_{ap^s}}{P_{ap^s}}
\]
is a well-defined element of $\Qp$ with $\vp(\widehat{c}_p(a))\in\{-2,-1,0\}$;
it is \emph{not} rational (rational reconstruction sits at the noise floor) and
it \emph{depends on $a$} (different $a$ agree to $0$ digits).  The renormalised
$p$-adic defect is thus a function on the tower tree, not a constant.
\end{finding}

\subsection{The archimedean contrast: positivity versus grading}\label{ss:contrast}

Archimedeanly the same two linear forms ride the \emph{same} ray: $\log|\Ip_n|/n$
and $\log|\Ih_n|/n$ both tend to $\lambda_2=-2.47237\dots$, so
$c=\lim_n\Ih_n/\Ip_n$ is a finite nonzero constant, namely
$c=-1/(2\zeta(2))=-3/\pi^2=-0.30396\dots$ \PROVED\ in the companion
\cite{PaperB-inprep}; and that constant is exactly what forces the
Brown--Zudilin class off the minimal ray, because a positivity relation
expressing $\Ip_n$ as a positive combination of two further cell integrals pins
the cellular class to the ray $\rho_B$, of which the impurity
$\zeta(5)+2\zeta(2)\zeta(3)$ of the top period is the visible shadow
\cite{PaperB-inprep}.  $p$-adically that degeneracy is
broken: the Frobenius grading $\mathrm{diag}(1,p^3,p^5)$ puts $Q$, $\Ph$, $P$
into three \emph{distinct} valuation strata $0$, $-3L$, $-5L$
(Finding~\ref{find:strata}), so the two forms cannot ride a common ray and the
defect collapses to $c_p=0$ (Finding~\ref{find:cp}).  The strategic asymmetry
this exposes is worth stating once, plainly: \emph{the $p$-adic side of this
family is structurally clean but arithmetically weak, while the archimedean side
is arithmetically strong but structurally obstructed}; there is no purity-defect
obstruction $p$-adically at all, and what remains there is purely a size budget
(\S\ref{sec:open} and Appendix~\ref{app:LSZ}).

% =====================================================================
\section{The exclusion theorem}\label{sec:exclusion}
% =====================================================================

We now come to the negative result, which we regard as the main arithmetic
content of the paper.  A natural and frequently voiced expectation is that,
since $b_n/a_n\to\zeta(3)$, the $p$-adic tower limit $\Lambda_a$ ought to be a
$\zeta_p(3)$-avatar --- ``one ratio, two completions''.  It is not.

\subsection{The reference values}\label{ss:refvals}

We use $\zeta_p(s)=L_p(s,\omega^{1-s})$ as in \S\ref{ss:notation}, computed from
a Volkenborn-integral series derived from scratch and validated on $980$
independent interpolation identities; the construction, the proof that the
series \emph{is} $L_p$, and the full validation record are in
Appendix~\ref{app:KL}.  We record here only the valuations, which matter for the
strata argument:
\[
  \vp(\zeta_p(3))=0\ \ (p\ge5),\qquad
  \vp(\zeta_p(5))=\begin{cases}-1,&p=5,\\ 0,&p=7,11,13,\end{cases}
\]
the $p=5$ case being the one where $\omega^{1-5}=\omega^{-4}$ is trivial because
$(p-1)\mid4$.  \VER{$p\in\{5,7,11,13\}$, $12$--$22$ digits.}

\subsection{No low-height $\QQ$-affine relation}\label{ss:LLL}

The search instrument is an exact-Fraction LLL with the residue column weighted
by $p^{K}$, so that only genuine relations survive; it was validated first on
planted relations and on true relations among Iwasawa logarithms
($\log_p4-2\log_p2=0$ at height $2$; $\log_p6-\log_p2-\log_p3=0$ at height $1$;
a planted $5x-7\zeta_p(3)-3=0$ at height $7$) --- \VER{recovered at every prime
tested}.  For an $m$-term relation certified to $K$ digits the LLL noise floor
is $p^{K/m}$; a genuine relation must appear \emph{below} the floor.

\begin{finding}[Exclusion]\label{find:exclusion}
\VER{$p=5,7,11,13$; $K=12$--$27$ certified digits; $a=1,\dots,12$ with $p\nmid a$}
No relation
\[
  c_0+c_1\zeta_p(3)+c_2\Lambda=0
  \qquad\text{or}\qquad
  c_0+c_1\zeta_p(5)+c_2\Lambda=0
\]
with integer height $\lesssim10^4$ holds for $\Lambda$ any of $\Lambda_a$
(Ap\'ery), $\Lap_a$, $\Lah_a$ (Brown--Zudilin).  Every candidate returned sits
\emph{at} the noise floor, i.e.\ is nothing.  The same is true of:
the $2$-term test $(\zeta_p(3),\Lambda_1)$; the normalisation-free test
$\Lambda_a/\Lambda_1\in\QQ$; the tests against $\log_p2$, $\log_p3$, and the
unit root of the weight-$4$ level-$8$ newform $\eta(2z)^4\eta(4z)^4$ that
Beukers attaches to the Ap\'ery numbers \cite{Beukers1987}; the $4$-term
test $(1,\zeta_p(3),\log_p2,\Lambda)$; and algebraicity of degree $\le3$.
Tables~\ref{tab:lll1} and~\ref{tab:lll2} record the heights against the floors.
\end{finding}

\begin{table}[h]
\centering\small
\caption{Ap\'ery pair: best LLL height for $(1,\zeta_p(3),\Lambda_1)$ against
the $3$-term noise floor $p^{K/3}$.}
\label{tab:lll1}
\begin{tabular}{@{}cccc@{}}
\toprule
$p$ & $K$ (certified digits) & best height & noise floor $p^{K/3}$\\
\midrule
5 & 23 & $1.0\times10^5$ & $2.3\times10^5$\\
7 & 21 & $4.9\times10^5$ & $8.2\times10^5$\\
11 & 15 & $4.4\times10^4$ & $1.6\times10^5$\\
13 & 15 & $1.9\times10^5$ & $3.7\times10^5$\\
\bottomrule
\end{tabular}
\end{table}

\begin{table}[h]
\centering\small
\caption{The wider hunt at $p=5$, $K=27$: best LLL height found for each
(target, basis) pair.  Nothing anywhere below the floor.  $C_1:=\Lambda_1-h_p(1)$.}
\label{tab:lll2}
\begin{tabular}{@{}lcccccc@{}}
\toprule
target & $(1,\zeta_p(3),\cdot)$ & $(1,\zeta_p(5),\cdot)$ & $(1,\log_p2,\cdot)$
       & $(1,\log_p3,\cdot)$ & $(1,\text{unit root},\cdot)$
       & $(1,\zeta_p(3),\log_p2,\cdot)$\\
\midrule
$\Lambda_1$ & 1.2e6 & 1.3e6 & 1.5e6 & 9.6e5 & 8.0e5 & 3.5e4\\
$\Atil_1$ & 1.1e6 & 1.1e6 & 1.2e6 & 1.2e6 & 9.6e5 & 2.1e4\\
$C_1$ & 1.3e6 & 1.7e6 & 6.7e5 & 7.4e5 & 8.2e5 & 2.5e4\\
$h_p(1)$ & 6.3e5 & 1.0e6 & 1.0e6 & 6.8e5 & 1.2e6 & 4.8e4\\
\midrule
\textbf{floor} & \textbf{2.0e6} & 2.0e6 & 2.0e6 & 2.0e6 & 2.0e6 & \textbf{5.2e4}\\
\bottomrule
\end{tabular}
\end{table}

The Brown--Zudilin rows behave identically: best heights $6\times10^3$ to
$1.4\times10^4$ for $(1,\zeta_p(5),\Lap_a)$ and $(1,\zeta_p(3),\Lah_a)$,
$a=1,2$, at $K=12$--$18$, against floors $1.5\times10^4$ to $1.7\times10^4$.
Nothing below the floor.  Also excluded: $h_p(1)\in\QQ$ (height
$5.2\times10^{10}$ vs floor $6.8\times10^{10}$) and $h_p(1)$ algebraic of
degree $2$.

\begin{remark}[why the margins are already decisive]\label{rem:BVheights}
Every Beukers--Vlasenko-style Frobenius constant --- $-\zeta_p(3)/3$,
$-8\zeta_p(3)/25$, $-35\zeta_p(3)/108$, $-\zeta_p(5)/5$, and their companions
\cite{BV2025} --- has height $\le10^2$.  These are excluded by
Finding~\ref{find:exclusion} with a margin of two to four orders of magnitude.
So the negative is not a close call for the candidates one would actually
propose.
\end{remark}

\subsection{The sharp form of the negative}\label{ss:sharp}

The LLL evidence bounds heights.  There is a stronger and entirely structural
statement which shows that no sharpening of estimates can help.

\begin{finding}[polynomial divergence]\label{find:diverge}
\VER{all $n\le5000$; $p\in\{5,7,11,13\}$}
\begin{gather*}
  \vp\bigl(b_n-\zeta_p(3)a_n\bigr)-\vp(a_n)=-3\floor{\log_pn}+O(1),\\
  \vp\bigl(Q_n\zeta_p(5)-P_n\bigr)-\vp(Q_n)=-5\floor{\log_pn}+O(1).
\end{gather*}
Equivalently
$|a_n\zeta_p(3)-b_n|_p\asymp n^3|a_n|_p$ and
$|Q_n\zeta_p(5)-P_n|_p\asymp n^5|Q_n|_p$: the $p$-adic linear forms
\emph{diverge polynomially} instead of converging.
\end{finding}

A Bel-type $p$-adic irrationality criterion \cite[Lemma~4]{LSZ2025} requires
$|a_n+b_n\xi|_p\le\exp(-\alpha n)$.  Ours grows like a power of $n$.  The gap is
therefore infinite, not finite, and no sharpening of any estimate can rescue
these pairs $p$-adically.  Note the logical relation to
Theorem~\ref{thm:exists}: the normalisation by $p^{ws}$ is precisely the act of
dividing this divergence out, and what survives, $\Lambda_a$, is then no longer
a linear form in anything.

\subsection{The structural reason: order $3$ is too short}\label{ss:order3}

Beukers--Vlasenko \cite[Prop.~1.3]{BV2025} write the Frobenius structure of a
MUM operator as $\mathcal{A}\bigl(y_i(t^p)\bigr)=p^{i}\sum_j\alpha_jy_{i-j}(t)$,
so an operator of order $3$ has only $\alpha_1$ and $\alpha_2$.  Their
Theorem~1.5 evaluates these:
\begin{equation}\label{eq:BVgen}
  \alpha_j=[x^j]\,\Gp(x)e^{-\Gp'(0)x}
         =[x^j]\exp\Bigl(-\sum_{m\ge2}\zeta_p(m)\frac{x^m}{m}\Bigr).
\end{equation}
Because there is no $\zeta_p(1)$ and $\zeta_p(2)=0$, this gives
\[
  \alpha_1=\alpha_2=0,\qquad
  \alpha_3=-\tfrac{\zeta_p(3)}{3},\quad
  \alpha_5=-\tfrac{\zeta_p(5)}{5},\quad
  \alpha_6=\tfrac{\zeta_p(3)^2}{18},\ \dots
\]
The Ap\'ery $\zeta(3)$ operator
$\theta^3-t(34\theta^3+51\theta^2+27\theta+5)+t^2(\theta+1)^3$ has order $3$.
Its $p$-adic Frobenius structure is therefore \emph{trivial}: the available
slots are literally zero, and $\zeta_p(3)$ first appears at $\alpha_3$, which
needs order $\ge4$.  In particular $\Lambda$ is definitely not an $\alpha_j$.
(That the Ap\'ery operator is a symmetric square of an order-$2$ operator,
leaving no room whatever, is \RECALLED.)

\subsection{Consistency with the literature}\label{ss:consistency}

Finding~\ref{find:exclusion} and Finding~\ref{find:diverge} are consistent with,
and arguably explain, two facts in the literature.  First, Calegari's $p$-adic
irrationality theorem for $\zeta_p(3)$ at $p=2,3$ \cite{Calegari2005} does
\emph{not} use the Ap\'ery pair: it uses overconvergent $p$-adic modular forms.
Second, the successful $p$-adic constructions of Lai--Sprang--Zudilin
\cite{LSZ2025} obtain their linear forms from a Volkenborn integral of a
rational function of degree $-3$, exploiting two vanishing mechanisms
unavailable archimedeanly: purity by degree
($\sum_kr_{n,1,k}=\lim_{t\to\infty}tR_n(t)=0$, which kills the $\zeta_2(3)$
coefficient outright and leaves a genuine two-term form), and $\zeta_p(2k)=0$,
which kills the even-weight terms for free.  Neither mechanism is available to
the Ap\'ery or Brown--Zudilin pairs, whose $p$-adic content, as we now show, is
not a linear form at all but a $\Gp$-value.

%% FIXED (MAJOR M-2, evidence-class): this was the only theorem-class environment in the paper
%% FIXED carrying \VER and no \PROVED, in violation of \S\ref{ss:notation}'s own convention. It is
%% FIXED now a Finding, and the unsupported quantifier "at any prime" has been replaced by the
%% FIXED range actually swept (p in {5,7,11,13}, heights <~ 10^4, n <= 5000), with the
%% FIXED extrapolation stated as such.
\begin{finding}\label{cor:noroute}
\VER{as in Findings~\ref{find:exclusion} and \ref{find:diverge}}
Within the range swept --- $p\in\{5,7,11,13\}$, relation heights $\lesssim10^{4}$,
$n\le5000$ --- there is no route to the irrationality of any $\zeta_p(2k+1)$ through the
classical Ap\'ery pair or the Brown--Zudilin pair.
\end{finding}

We state the wider claim separately, as an expectation rather than a result, because it is the
one place in this paper where we extrapolate past the sweep.

\begin{remark}
We expect the same at every prime, and Finding~\ref{find:diverge} is the reason: the
obstruction found is not a height bound that a longer search might clear, but polynomial
divergence of the linear forms, a structural property of the pair that does not depend on $p$.
That is an argument, not a proof, and no numerical evidence outside the four primes above
supports it.
\end{remark}

% =====================================================================
\section{Identification: what the tower limits are}\label{sec:identify}
% =====================================================================

The exclusions of \S\ref{sec:exclusion} are not accidents of height.  In this
section we identify $\Lambda_a$ completely: it is an explicit, infinite,
convergent assembly of two ingredients, each of which has a closed form, and
each of which involves \emph{all} the odd $p$-adic zeta values at once.

\subsection{The harmonic split and $h_p(a)$}\label{ss:hp}

Ap\'ery's formula
\begin{gather*}
  b_n=\sum_kA(n,k)\bigl[H_3(n)+T(n,k)\bigr],\\
  A(n,k)=\bin{n}{k}^2\bin{n+k}{k}^2,\qquad
  T(n,k)=\sum_{m\le k}\frac{(-1)^{m-1}}{2m^3\bin{n}{m}\bin{n+m}{m}},
\end{gather*}
splits $f=b/a$ into $H_3+c_n/a_n$, where $H_3(n)=\sum_{k\le n}k^{-3}$.  The
harmonic part has an exact tower limit.

\begin{theorem}\label{thm:hp}
\PROVED.  For $p\nmid a$,
\[
  p^{3s}H_3(ap^s)=\sum_{i}p^{3i}S_i(a),
  \qquad S_i(a)=\sum_{k\le ap^{i},\ p\nmid k}k^{-3}
\]
(the sum over $i$ with $ap^i\ge1$; over $0\le i\le s$ when $1\le a<p$), and
consequently
\[
  h_p(a):=\lim_{s\to\infty}p^{3s}H_3(ap^s)
  =\sum_{i}p^{3i}S_i(a)
  =\sum_{\substack{x\in\ZZ[1/p]\\ 0<x\le a}}x^{-3}\ \in\ \Qp .
\]
The series converges: $\vp\bigl(p^{3i}S_i(a)\bigr)\ge3i$.
\end{theorem}

\begin{proof}
Splitting $H_3$ at $p\mid k$ gives
$H_3(n)=p^{-3}H_3(\floor{n/p})+\sum_{k\le n,\,p\nmid k}k^{-3}$; iterating and
multiplying by $p^{3s}$ telescopes to the displayed identity.  Re-indexing
$x=k/p^{i}$ identifies the total with the sum over $x\in\ZZ[1/p]$, $0<x\le a$,
of $x^{-3}$: each such $x$ is uniquely $k/p^i$ with $p\nmid k$, and
$\vp(x^{-3})=3i$, so the terms tend to $0$ $p$-adically and the rearrangement is
legitimate.
\end{proof}

\VER{$h_p(a)$ against $p^{3s}H_3(ap^s)$ for $a=1,2,3$, $s\le3$,
$p\in\{5,7,11,13\}$: agreement $4,9,14,19$ digits.}  Note that the harmonic
tower converges at $5$ digits per level at weight $3$ --- faster than the Ap\'ery
tower --- exactly as Table~\ref{tab:harmonicrate} predicts.  Only the terms with
$i<K/4$ are needed for $K$ digits, so $h_p$ is computable to $40$+ digits in
milliseconds; this is what makes it usable as a reference constant in
Table~\ref{tab:lll2}.

\subsection{Kazandzidis limits in closed form}\label{ss:kaz}

\begin{definition}\label{def:kaz}
For integers $A>B\ge1$ set $c_p(A,B):=\lim_{t\to\infty}\bin{Ap^t}{Bp^t}\in\Qp$.
The limit exists by Jacobsthal--Kazandzidis \cite{JacobsthalKazandzidis}, at
$3$ digits per level (Finding~\ref{find:kazrate}).
\end{definition}

\begin{theorem}\label{thm:kaz}
%% FIXED (MINOR): the \VER range overstated one cell per prime --- t7_kaz.py drops K to 6 when
%% FIXED a*p^5 >= 4e6, so the pair (p^2,p) was checked at 6 digits, not 12 (kaz_p11.txt:
%% FIXED "c_p(121,11) ... MATCH(all 6)"). The claim itself is correct; only the range is now
%% FIXED stated as the run actually performed it.
\PROVED\ (from Morita's factorial factorisation and \cite[eq.~(2)]{BV2025}), and
\VER{$12$ pairs $(A,B)$ $\times$ $p\in\{5,7,11\}$, all digits matching, zero
mismatches, at $K=12$ digits except for the pair $(p^2,p)$, which the run
truncates to $K=6$}.
Let $D=A-B$ and, for $l\ge1$,
$A_l=\floor{A/p^l}$, $B_l=\floor{B/p^l}$, $D_l=\floor{D/p^l}$.  Put
\[
  T(A,B):=\prod_{l\ge1}
    \frac{(-1)^{A_l+1}\Gp(A_l+1)}
         {(-1)^{B_l+1}\Gp(B_l+1)\cdot(-1)^{D_l+1}\Gp(D_l+1)}
\]
(a finite product, equal to $1$ when $A<p$).  Then
\begin{equation}\label{eq:kaz}
\begin{aligned}
  c_p(A,B)
  &= -\,p^{\vp\binom{A}{B}}\ T(A,B)\
      \frac{\Gp(A+1)}{\Gp(B+1)\,\Gp(D+1)}\ e^{\,\mathcal{Z}(A,B)},\\[2pt]
  \mathcal{Z}(A,B)
  &= -\!\!\sum_{\substack{m\ge3\\ m\ \mathrm{odd}}}\!\!
        \zeta_p(m)\bigl(A^m-B^m-D^m\bigr)\frac{p^m}{m(1-p^m)}.
\end{aligned}
\end{equation}
\end{theorem}

\begin{proof}[Proof sketch]
Morita's factorisation \cite{Morita1975}
$n!=p^{\vp(n!)}\prod_{i\ge0}(-1)^{n_i+1}\Gp(n_i+1)$, $n_i=\floor{n/p^i}$,
applied to $A p^t$, $Bp^t$, $Dp^t$ turns $\bin{Ap^t}{Bp^t}$ into
$p^{\vp\binom{A}{B}}$ times a ratio of $\Gp$-values at the arguments
$\floor{Ap^t/p^i}+1$ etc.  The arguments with $i\le t$ are
$Ap^{t-i}+1$, and $\Gp(1+x)=-\Gp(x)$ on $p\ZZ_p$ moves them onto $p$-power
points; the arguments with $i>t$ produce exactly the finite tail $T(A,B)$,
which is independent of $t$.  Applying \cite[eq.~(2)]{BV2025},
$\log\Gp(x)=\Gp'(0)x-\sum_{m\ge2}\zeta_p(m)x^m/m$, at $x=Ap^{k}$, $Bp^{k}$,
$Dp^{k}$ and summing the resulting geometric series over $k$ gives the
exponential factor; the $\Gp'(0)$ terms cancel because $A=B+D$, and the
even-$m$ terms vanish because $\zeta_p(2k)=0$.
\end{proof}

The verified pairs include cases with carries, with $p\mid A$, and with
$\vp(c_p)=1$:
$(A,B)=(2,1),(3,1),(3,2),(4,1),(5,2),(7,3),(p,1),(p,2),(p+1,1),(p+1,p),(2p,p),(p^2,p)$.

The special case that drives the Ap\'ery tower is worth displaying.

\begin{corollary}\label{cor:cp21}
\PROVED\ and \VER{to \emph{all} certified digits: $30$ at $p=5$ (from
$n=2\cdot5^9$), $24$ at $p=7$, $21$ at $p=11$, $18$ at $p=13$}
\[
  c_p:=\lim_{s\to\infty}\bin{2p^s}{p^s}
  =2\exp\Bigl(-\!\!\sum_{\substack{m\ge3\\ m\ \mathrm{odd}}}\!\!
    \zeta_p(m)\,\frac{(2^m-2)\,p^m}{m(1-p^m)}\Bigr).
\]
\end{corollary}
\begin{proof}
Specialise \eqref{eq:kaz} at $(A,B)=(2,1)$: $\vp\binom21=0$, $T(2,1)=1$, and
$\Gp(3)/\Gp(2)^2=-2$.
\end{proof}

\begin{remark}[the quantitative form of the exclusions]\label{rem:trunc}
\VER{$p=5,7,11,13$} Truncating the series in Corollary~\ref{cor:cp21} after the
$\zeta_p(3)$ term reproduces only $4$--$5$ correct digits out of $24$--$30$.
The constant genuinely needs the entire tower of odd $p$-adic zeta values.
This is the mechanism behind Finding~\ref{find:exclusion}, quantified.
\end{remark}

\subsection{The structure theorem}\label{ss:structure}

Along $n=p^s$ one has $\vp\bin{p^s}{k}=s-\vp(k)$, so writing $k=jp^{s-i}$ with
$p\nmid j$, $0<j\le p^{i}$, the Ap\'ery summand $A(n,k)$ has valuation $2i$ and
\[
  \lim_{s\to\infty}A\bigl(p^s,jp^{s-i}\bigr)
  =\bigl[c_p(p^{i},j)\,c_p(p^{i}+j,j)\bigr]^2 .
\]
The inner Ap\'ery weight at $m=j'p^{s-i'}$ contributes
\[
  \frac{p^{3s}}{2m^3\bin{n}{m}\bin{n+m}{m}}
  \ \longrightarrow\
  \frac{(-1)^{j'-1}p^{3i'}}{2\,j'^{3}\,c_p(p^{i'},j')\,c_p(p^{i'}+j',j')},
\]
of valuation $2i'$.  Assembling:

\begin{finding}[Structure theorem]\label{find:structure}
\VER{$p=5,7,11,13$, block by block; each predicted block confirmed and the
residual moving to exactly the next block}  With
$\Atil_a=\lim_sa_{ap^s}$ and $\Btil_a=\lim_sp^{3s}b_{ap^s}$ (so
$\Lambda_a=\Btil_a/\Atil_a$),
\begin{align}
  \Atil_1&=1+\sum_{i\ge0}\ \sum_{\substack{0<j\le p^{i}\\ p\nmid j}}
    \bigl[c_p(p^{i},j)c_p(p^{i}+j,j)\bigr]^2,
    \label{eq:Atil}\\
  \Lambda_1&=h_p(1)+\frac{1}{\Atil_1}
    \sum_{(i,j),(i',j')}\bigl(\text{block of valuation }2(i+i')\bigr),
    \label{eq:Lam}
\end{align}
where the leading term of \eqref{eq:Atil} is the $k=0$ term of the sum defining
$a_n$, the $(i,j)$ block of \eqref{eq:Atil} has valuation exactly $2i$, and the
leading blocks of \eqref{eq:Lam} are
\[
  (0,0):\ \frac{c_p(2,1)}{2},
  \qquad
  (0,1):\ c_p(2,1)^2\sum_{j=1}^{p-1}
     \frac{(-1)^{j-1}p^3}{2j^3\,c_p(p,j)\,c_p(p+j,j)} \ =:\ X_3 .
\]
\end{finding}

Table~\ref{tab:blocks} is the block-by-block verification: in each row the
right-hand side is the truncation of the block sum and the entry is the
valuation of the first discrepancy against the directly computed constant.

\begin{table}[h]
\centering\small
\caption{Block-by-block verification of Finding~\ref{find:structure}.}
\label{tab:blocks}
\begin{tabular}{@{}lll@{}}
\toprule
check & first discrepancy at $\vp$ & predicted\\
\midrule
$\Atil_1=1+c_p(2,1)^2$ & 3 & $\ge2$\\
$\Atil_1=1+c_p(2,1)^2+\sum_j[c_p(p,j)c_p(p+j,j)]^2$ & 6 & $\ge4$\\
$(\Lambda_1-h_p(1))\Atil_1=c_p(2,1)/2$ & 3 & $\ge2$\\
$(\Lambda_1-h_p(1))\Atil_1=c_p(2,1)/2+X_3$ & 4 ($p=5$), 5 ($p=11,13$), 6 ($p=7$) & $\ge4$\\
\bottomrule
\end{tabular}
\end{table}

%% FIXED (MINOR): the last row is a MISS against the generating script's own finer prediction
%% FIXED (struct_p5.txt: "discrepancy at v_p = 4 (expect 5)"); LAMBDA_HUNT:203 relaxed the column
%% FIXED to ">= 4" and the paper inherited the relaxed version. Now stated.
\begin{remark}
Every row confirms the predicted block, in the sense that the residual moves to the valuation of
the next block.  One row is worth flagging: for the last check the generating script predicted
the first discrepancy at $\vp=5$ and observed it at $\vp=4$ at $p=5$
(\texttt{struct\_p5.txt}).  The block decomposition predicts only $\ge4$ there, which is what
the table records and what holds; but the sharper $\vp=5$ that one might expect from the
pattern of the other primes does \emph{not} hold at $p=5$.
\end{remark}

\begin{remark}\label{rem:complete-not-finite}
The identification is complete in the sense that every ingredient of
$\Lambda_a$ is named and evaluated in closed form.  What is \emph{not} finite is
the number of ingredients: each new level of the tower contributes genuinely new
Kazandzidis limits $c_p(p^i,j)$.  This is the precise sense in which
$\Lambda_a$ is ``not a new constant'' and simultaneously ``not a known one''.
\end{remark}

\subsection{Why every exclusion holds}\label{ss:whyexcl}

Combining Theorem~\ref{thm:kaz} with Finding~\ref{find:structure}: the tower
limits live in the $\Gp$-value ring.  Each is a convergent $\QQ$-combination of
products of exponentials of infinite $\QQ$-series in
$\zeta_p(3),\zeta_p(5),\zeta_p(7),\dots$ with $p$-power-weighted rational
coefficients.  No finite low-height $\QQ$-relation involving $\zeta_p(3)$ alone
--- or $\zeta_p(5)$ alone --- can hold, and Remark~\ref{rem:trunc} quantifies
how badly a $\zeta_p(3)$-only model fails ($5$ digits out of $24$ for
$\Lambda_1$; $4$--$5$ out of $24$--$30$ for $c_p$).  The exclusion table is
therefore not a collection of near-misses but the expected shadow of a
structure theorem.

\subsection{Positioning: values versus Taylor coefficients}\label{ss:BVposition}

Equation~\eqref{eq:BVgen} and equation~\eqref{eq:kaz} involve the same function,
\[
  G(x):=\Gp(x)e^{-\Gp'(0)x}=\exp\Bigl(-\sum_{m\ge2}\zeta_p(m)\frac{x^m}{m}\Bigr),
\]
and they extract different information from it.  Beukers--Vlasenko take the
\emph{Taylor coefficients of $G$ at $0$}: these are the Frobenius structure
constants $\alpha_j$, and for an order-$3$ operator only $\alpha_1,\alpha_2$
exist, both zero.  Our tower limits are \emph{values of $G$} at the $p$-adic
points $x=Ap^{k}$, multiplied together over $k$ and assembled over the tower
tree.  That is the precise sense in which the constants of this paper are
$p$-adic Frobenius data, and simultaneously the precise reason they are not
$\QQ$-affine in any single $\zeta_p(m)$.

Three further points of orientation.

\begin{enumerate}[label=(\arabic*),leftmargin=2em]
\item \emph{Terminology.}  The Bloch--Vlasenko / Roy--Vlasenko \emph{Frobenius
  constants} $\kappa_n$ \cite{RoyVlasenko,Kerr} are archimedean periods
  (expansion coefficients of a generalised gamma function
  $\kappa(s)=s^2\int_0^1t^{s-1}\delta(t)\,dt$), conjecturally $\QQ$-combinations
  of products of zeta values of total weight $n$.  The Beukers--Vlasenko
  \emph{Frobenius structure constants} $\alpha_j$ \cite{BV2025} are the
  $p$-adic objects of \eqref{eq:BVgen}.  Our $\Lambda_a$ is the $p$-adic
  analogue of neither: it is the analogue of the Ap\'ery limit itself, i.e.\ of
  a $\kappa$-type period.
\item \emph{Why the naive expectation fails.}  The archimedean Ap\'ery limit is
  a value at one singular point.  The $p$-adic ``limit'' is a sum over the whole
  $p$-adic tree, and it therefore collects $G$ at \emph{every} level.  The
  expectation $\Lambda\leftrightarrow\zeta_p(3)/6$ compares an object of the
  first kind with an object of the second.
\item \emph{Adjacent literature.}  We also compared against the Frobenius-constant
  line of \cite{RoyVlasenko,Kerr} and the interpolated-Ap\'ery-number and
  quasiperiod circle of \cite{GolyshevZagier}; the objects there are
  archimedean $\kappa$-type periods and $q$-expansion data, and none of them is
  our $\Lambda$.
\item \emph{Novelty.}  We found no source naming these limits, and none stating
  the closed form of Corollary~\ref{cor:cp21}.  The ingredients --- Morita's
  factorisation, Jacobsthal--Kazandzidis, \cite[eq.~(2)]{BV2025} --- are all
  classical, so \eqref{eq:kaz} is presumably folklore-reachable; the derivation
  above is in any case self-contained and the formula is verified independently.
  \VER{as in Theorem~\ref{thm:kaz}.}
\end{enumerate}

\subsection{The one-generator refutation}\label{ss:onegen}

It is natural to hope that the family $\{\Lambda_a\}_{p\nmid a}$ is generated
over $\QQ$ by a single constant $\omega_p$ --- that the tower base $a$ enters
only through the elementary data $(a_a,b_a)$.  It does not.

\begin{finding}\label{find:onegen}
%% FIXED (CRITICAL knock-on): ``handed over for free by Prop 2.5(iv)'' is not universal ---
%% FIXED (iv) needs the p-unit-cell hypothesis. The floor of 3 is here a MEASURED residual, which
%% FIXED is what the table reports, and Prop 2.5(iv) explains it on the p-unit cells only.
\VER{$p=5,7,11,13$; $a\in\{1,\dots,12\}$, $p\nmid a$; $12$--$27$ certified
digits}  Every ansatz in Table~\ref{tab:ansatz} fails, each at the measured
trivial floor of $3$ residual digits.  On the $p$-unit cells that floor is
exactly what Proposition~\ref{prop:invariance}(iv) hands over for free, and by
Remark~\ref{rem:ivsharp} no more than that is free; on the remaining cells it is
an empirical observation only.  A closing ansatz would show residual
agreement $\approx$ available digits.
\end{finding}

\begin{table}[h]
\centering\small
\caption{The one-generator hunt.  Entries are residual digits of agreement out
of $12$--$18$ available; the trivial floor is $3$.}
\label{tab:ansatz}
\begin{tabular}{@{}ll@{}}
\toprule
ansatz & result\\
\midrule
$\Lambda_a=\alpha b_a+\beta a_a$ & residual $0$--$1$ --- fails outright\\
$\Lambda_a=\alpha f(a)+\beta$ & residual $3$ --- trivial floor only\\
$\Lambda_a=\mathrm{const}$; $\Lambda_a=\alpha f(a)$ & $0$ / $3$ --- fails\\
M\"obius $(\alpha b_a+\beta a_a)/(\gamma b_a+\delta a_a)$ & residual $0$--$2$ --- fails\\
M\"obius in $(a_a,b_a,1)$ & residual $0$--$2$ --- fails\\
M\"obius in the state vector $(a_a,b_a,a_{a+1},b_{a+1})$ & residual $-1$--$0$ --- fails\\
$\Atil_a=\alpha a_a+\beta b_a$ (and variants) & residual $3$ --- trivial floor\\
$\Btil_a=\alpha a_a+\beta b_a$ & residual $3$--$4$ --- trivial floor\\
$\Atil$ satisfies the Ap\'ery recurrence? & residual $\vp=3$ at every $n$ --- no\\
$\delta_a/\delta_b\in\QQ$, $\delta_a:=\Lambda_a-f(a)$ & all reconstructions at the noise floor\\
$C_a/C_b\in\QQ$, $C_a:=\Lambda_a-h_p(a)$ & all reconstructions at the noise floor\\
$(1,\Lambda_a,\Lambda_b)$ spans $2$ over $\QQ$ & heights $9$e$4$--$6$e$5$ vs floors $1$--$7$e$5$ --- noise\\
$\Lambda_1,\Atil_1,\Btil_1,C_1,h_p(1)$ algebraic of degree $2$ or $3$ & all at the noise floor\\
\bottomrule
\end{tabular}
\end{table}

Finding~\ref{find:structure} explains why no ansatz can close.  $\Lambda_a$ is a
sum over a two-index family of blocks of valuation $2(i+i')$, and each new level
contributes genuinely new Kazandzidis limits.  There is no finite-dimensional
$\QQ$-structure to find, hence no $\omega_p$.

One classical-looking candidate deserves an explicit burial.  The Casoratian of
the Ap\'ery pair is
\begin{equation}\label{eq:casoratian}
  a_{n+1}b_n-a_nb_{n+1}=-\frac{6}{(n+1)^3},
\end{equation}
\PROVED\ from \eqref{eq:apery-rec}: $W_n:=a_{n+1}b_n-a_nb_{n+1}$ satisfies
$(n+1)^3W_n=n^3W_{n-1}$ with $W_0=-6$.  It is a pure $p$-power,
$\vp(W_n)=-3\vp(n+1)$, so it carries \emph{no} normalising information about a
tower: it is invisible at $n=ap^s$ with $p\nmid a$.  This is exactly why the
Casoratian-corrected ans\"atze collapse onto the plain ones and fail with them.

\begin{remark}[a methodological trap, recorded]\label{rem:trap}
The Fermat quotients $q_p(2),q_p(3),q_p(5)$, the Wolstenholme quotients
$H_{p-1}/p^2$ and $H^{(2)}_{p-1}/p$, and $B_{p-3}$ are \emph{rational numbers}
(e.g.\ $q_5(2)=3$, $H_6/49=1/20$, $B_{10}=5/66$).  Placing them in an
integer-relation search against $1$ is vacuous: it ``finds'' relations such as
$q_p(2)-9=0$ and reports false candidates.  They are congruence data modulo
small powers of $p$, not independent $p$-adic constants.  Similarly
$\Gp(1/2),\Gp(1/3),\Gp(1/4)$ are algebraic --- $\Gp(1/2)^2=\pm1$ by the
reflection formula, and the others are Gauss-sum values by Gross--Koblitz
\cite{GrossKoblitz1979} --- so they are already covered by the
algebraicity tests of Table~\ref{tab:ansatz}.  The $\Gp$ that matters here is
$\Gp$ at $p$-power arguments, which is exactly what Theorem~\ref{thm:kaz}
produces.
\end{remark}

% =====================================================================
\section{Existence without archimedean limits}\label{sec:noarch}
% =====================================================================

We close the positive part of the paper with what we regard as the cleanest
single piece of evidence that the $p$-adic tower limit is an object about
Frobenius rather than about archimedean asymptotics.

Three of the fifteen sporadic sequences have \emph{no} archimedean Ap\'ery
limit.  For \eqref{eq:R2} the characteristic roots solve
$\lambda^2-a\lambda+c=0$ and for \eqref{eq:R3} they solve
$\lambda^2-2a\lambda+c=0$; the discriminants are
\[
  B:\ -27,\qquad (\delta):\ -128,\qquad (\eta):\ -16,
\]
complex-conjugate pairs of equal modulus.  Poincar\'e--Perron therefore supplies
no dominant solution and $B(n)/A(n)$ does not converge --- \VER{$0$ digits of
agreement between $n=590$ and $n=600$, against $380$ digits for the convergent
families}.  Every other sporadic sequence has real distinct roots.

%% FIXED (MINOR): the \VER range said p=5,7,11 for all three, but Finding \ref{find:twist}
%% FIXED records that (eta) degenerates entirely at p=5 (chi_5(5)=0). Range now split.
\begin{finding}\label{find:noarch}
\VER{$N=20\,000$; $a=1,2$; $p=5,7,11$ for $B$ and $(\delta)$, $p=7,11$ for $(\eta)$}
All three of $B$, $(\delta)$, $(\eta)$
have $p$-adic tower limits at every unramified $p\ge5$:
\begin{itemize}[leftmargin=1.5em]
\item $(\delta)$ outright, at rate $3$ (its character is trivial);
\item $B$ and $(\eta)$ after the twist, i.e.\ for
  $\Lambda^\chi_a=\lim_s\chi(p)^{s}p^{ws}B(ap^s)/A(ap^s)$, at rates $2$ and $3$
  respectively.
\end{itemize}
All three satisfy $(\mathrm{LB}_w^\chi)$ with a flat floor, with characters
$\chi_{-3}$, trivial, $\chi_5$ respectively.  The one exclusion is the ramified prime:
$(\eta)$ has $\chi_5(5)=0$ and no tower limit at $p=5$, by
Finding~\textup{\ref{find:twist}}.
\end{finding}

\begin{finding}[the twist is necessary, and the ramified case degenerates]\label{find:twist}
\VER{$p=5,7,11$; $E$ and $(\eta)$ at $p=7$; $B$, $C$, $F$, $(\zeta)$, $s_{18}$
at $p=5,11$}  For a $\chi$-twisted family at a prime with $\chi(p)=-1$
the untwisted tower does not converge at all ($0$ digits gained per level);
it converges after normalising by $\chi(p)^{-s}$.  This is the normalisation
demanded by the twisted master law $(\mathrm{LB}_w^\chi)$.  At a prime dividing
the conductor --- $(\eta)$ at $p=5$, where $\chi_5(5)=0$ --- the tower
degenerates entirely: $\vp(\Lambda)$ drifts and there is no convergence.  This
is the only ramified prime $\ge5$ among the fifteen families.
\end{finding}

The ramified case is benign at the level of the congruence, if not of the limit:
there $(\mathrm{LB}_w^\chi)$ reads $p^{w}B_nA_q\equiv0\pmod{p^{w}}$, which is
exactly the assertion that $B_n$ is $p$-integral, and \VER{$n\le500$}
$\min_n v_5(B_n)=0$ --- the pole simply does not develop at the ramified prime.

\begin{finding}[no cross-family relations]\label{find:cross}
\VER{$p=7$, all pairs of the fifteen families, $12$--$21$ digits}  Pairwise
integer-relation searches for $(1,\Lambda^{\mathrm{fam}_1}_1,
\Lambda^{\mathrm{fam}_2}_1)$ find nothing below the noise floor.  The tower
limits of different families all lie in the same ring --- the $\Gp$-value ring
of \S\ref{ss:whyexcl} --- but are not $\QQ$-linearly related to one another at
accessible heights.
\end{finding}

% =====================================================================
\section{Open problems}\label{sec:open}
% =====================================================================

\begin{problem}[Name the constants]\label{prob:name}
Is $h_p(1)=\sum_{x\in\ZZ[1/p],\,0<x\le1}x^{-3}$ irrational?  Transcendental
over $\QQ$?  We have excluded $h_p(1)\in\QQ$ at height $6.8\times10^{10}$ and
algebraicity of degree $2$ (\S\ref{ss:LLL}), but nothing is proved.  The same
question for $c_p=\lim_s\bin{2p^s}{p^s}$, for which Corollary~\ref{cor:cp21}
gives a closed form entirely in terms of $\zeta_p(3),\zeta_p(5),\dots$: is
$c_p$ transcendental?  Is it in the field generated by the odd $p$-adic zeta
values, and if so, is that containment strict?
\end{problem}

\begin{problem}[Upgrade the Findings]\label{prob:upgrade}
Prove the multi-digit master descent $(\mathrm{LB}_w^\chi)$ modulo $p^{w}$;
Paper~A \cite{PaperA-inprep} proves only the single-digit form modulo $p$.
Prove the scaled weight-$5$ law \eqref{eq:scaled5} with its depth $3s$.  Prove
the strata bounds of Finding~\ref{find:strata} with an explicit uniform $O(1)$,
and deduce $c_p=0$ (Finding~\ref{find:cp}) unconditionally.
\end{problem}

\begin{problem}[The rate law beyond $w=3$]\label{prob:rate}
Finding~\ref{find:harmrate} is verified only at $w=2,3$.  Prove
$r(\chi,w)=w$ for $\chi\ne1$ and $r(1,w)=w+1$ ($w$ even), $w+2$ ($w$ odd), and
prove the cap: that a tower limit assembled from Kazandzidis limits cannot gain
more than $3$ digits per level.  Explain the upward deviations at
$(\varepsilon)$, $(\zeta)$, $C$, $s_{18}$.
\end{problem}

\begin{problem}[Order $\ge4$: is there a $\zeta_p$-carrying Ap\'ery family?]\label{prob:order4}
By \S\ref{ss:order3} the value $\zeta_p(3)$ first occupies the Frobenius slot
$\alpha_3$, which requires an operator of order $\ge4$.  Is there an
Ap\'ery-like family of order $\ge4$ whose tower limit \emph{is}
$\QQ$-affine in $\zeta_p(3)$ --- i.e.\ for which the naive ``one ratio, two
completions'' expectation is true?  Our exclusion is a statement about the
order-$3$ and rank-$3$ families we can compute with, not about the phenomenon in
general.
\end{problem}

\begin{problem}[The intrinsic meaning of the twist]\label{prob:twist}
The factor $\chi(p)$ enters through Lemma~\ref{lem:K}, on the non-unit Frobenius
eigenvalue only; it is invisible in the $A$-row.  Any crystal-theoretic account
of $(\mathrm{LB}_w^\chi)$ must carry the information that for seven of the
fifteen families the second Frobenius eigenvalue is $\chi(p)p^{w}$ and not
$p^{w}$ --- i.e.\ that the family has a CM or quadratic-twisted modular
parametrisation.  What is the correct statement, and what happens at
$p\mid\cond(\chi)$, where the tower degenerates
(Finding~\ref{find:twist})?
\end{problem}

\begin{problem}[The $O(1)\to\delta n$ bridge]\label{prob:bridge}
This is, in our view, the field's question, and we state it as such.  No
individual $\zeta_p(2k+1)$ is known to be irrational for a single prime
$p\ge5$; the three known individual results all live at $p=2,3$.  Fitting the
same ledger shape to those three results (Appendix~\ref{app:LSZ}) gives
$\mathrm{margin}(p,w)=(w+3)\log p/(p-1)-w$, reproducing all three published
irrationality measures exactly \FIT, and predicts the two nearest misses:
\[
  \zeta_5(3):\ \text{deficit }0.586,\qquad
  \zeta_3(5):\ \text{deficit }0.606 .
\]
Closing either requires a denominator saving of $e^{\delta n}$ with
$\delta\approx0.59$ resp.\ $0.61$.  The descent congruences
$(\mathrm{LB}_w^\chi)$ are the single-prime shadow of exactly this kind of
saving --- but they save $O(1)$ per prime (bounded depth, exactly $w$), hence
$o(n)$ in total, where one needs $\delta n$.  \emph{Bridging $O(1)\to\delta n$ is
the whole problem.}  Appendix~\ref{app:LSZ} shows, by exact reproduction of the
Lai--Sprang--Zudilin ledger, that the $d_n^5$ budget there is tight: there is no
free denominator saving hiding in the present construction, so the saving must
come from a group action on a multi-parameter family
\cite{RhinViola} rather than from a sharper estimate on a fixed one.  We record, as the
complementary judgement, that we do \emph{not} regard linear independence of
$1,\zeta_p(3),\zeta_p(5)$ for a specific $p$ as the promising target: two-term
forms are already at the edge, and a three-term form would need the elimination
technology of \cite{FSZ2018,LaiSprang2023}, whose $p$-adic versions give
asymptotic counts rather than a fixed small index set.
\end{problem}

\begin{problem}[Is the $\Gp$-value ring the right home?]\label{prob:ring}
Finding~\ref{find:cross} says the tower limits of different families are not
$\QQ$-linearly related at accessible heights, while
Finding~\ref{find:structure} says they all lie in the ring generated by values
of $G(x)=\Gp(x)e^{-\Gp'(0)x}$ at $p$-power arguments.  Give that ring a
structure theory.  Is there a $p$-adic analogue of the weight filtration under
which $h_p$, $c_p(A,B)$ and $\Lambda_a$ have well-defined weights?
\end{problem}

% =====================================================================
\appendix
% =====================================================================

\section{The Kubota--Leopoldt implementation, and its validation}\label{app:KL}

Every exclusion in \S\ref{sec:exclusion} is only as good as the reference values
$\zeta_p(3)$, $\zeta_p(5)$.  We therefore built the Kubota--Leopoldt
$L$-function from scratch, proved that what we built is $L_p$, and validated the
implementation against an independently coded classical construction.  This
appendix is the credibility of the negative result.

\subsection{The series}\label{app:KLseries}

Start from $\zeta_p(s)=L_p(s,\omega^{1-s})$ \cite[Def.~8]{LSZ2025}, the
Volkenborn integral $\int_{\Zp}t^k\,dt=B_k$ with the convention $B_1=-1/2$, and
the expansion $\langle t+x\rangle^{1-s}=\langle x\rangle^{1-s}(1+t/x)^{1-s}$
valid for $|x|_p>1$.  These give
\[
  \zeta_p(s,x)=\frac{1}{s-1}\langle x\rangle^{1-s}
    \sum_{k\ge0}\binom{1-s}{k}B_kx^{-k},
\]
and hence, with $M$ any common multiple of $q_p$ and $\cond(\chi)$,
\begin{equation}\label{eq:KL}
  L_p(s,\chi)=\frac{1}{(s-1)M}
    \sum_{\substack{0<j<M\\ p\nmid j}}\chi(j)\langle j\rangle^{1-s}
    \sum_{k\ge0}\binom{1-s}{k}B_k\Bigl(\frac{M}{j}\Bigr)^{k}.
\end{equation}
Convergence is exact and effective: $\vp(B_k)\ge-1$ by von Staudt--Clausen and
$\vp\bigl((M/j)^k\bigr)=k\,\vp(M)$, so the $k$-th term has
$\vp\ge k\,\vp(M)-1$.

\begin{proposition}\label{prop:isKL}
\PROVED.  The right-hand side of \eqref{eq:KL} equals $L_p(s,\chi)$ for all
$s\in\Zp\setminus\{1\}$.
\end{proposition}

\begin{proof}
Set $s=1-n$ with $n\ge1$.  The inner sum truncates, and using
$B_n(x)=\sum_k\binom{n}{k}B_kx^{n-k}$ together with the classical
$B_{n,\psi}=F^{n-1}\sum_{a=1}^{F}\psi(a)B_n(a/F)$, the formula collapses to
\[
  L_p(1-n,\chi)=-\bigl(1-\chi\omega^{-n}(p)p^{n-1}\bigr)
    \frac{B_{n,\chi\omega^{-n}}}{n},
\]
which is the defining Kubota--Leopoldt interpolation property
\cite{KubotaLeopoldt1964}, for \emph{every} $n\ge1$.  The right-hand
series of \eqref{eq:KL} is continuous in $s\in\Zp$ (the convergence is uniform),
and $\{1-n:n\ge1\}$ is dense in $\Zp$, so the two agree on all of
$\Zp\setminus\{1\}$.
\end{proof}

\subsection{Validation}\label{app:KLvalid}

\begin{finding}[$980/980$]\label{find:980}
\VER{$p\in\{2,3,5,7,11,13,17,19\}$; every character $\omega^e$ with
$0\le e<p-1$; every $n\in[2,15]$; modulus $p^{12}$; $980$ identities; zero
mismatches}  The interpolation identity
$L_p(1-n,\chi)=-(1-\chi\omega^{-n}(p)p^{n-1})B_{n,\chi\omega^{-n}}/n$ holds for
our implementation against \emph{independently coded} classical generalised
Bernoulli numbers $B_{n,\psi}$: for trivial $\psi$, $B_n$ with the Euler factor
$1-p^{n-1}$; for $\psi$ of conductor $p$,
$p^{n-1}\sum_{a=1}^{p-1}\psi(a)B_n(a/p)$.
\end{finding}

Two further checks:

\begin{itemize}[leftmargin=1.5em]
\item \emph{Independence of the auxiliary modulus.}  \VER{$M=q_p\cdot m$,
  $m=1,\dots,5$} identical values in every case.  The formula \eqref{eq:KL}
  depends on a choice of $M$; the computed value does not.
\item \emph{A second, independent route.}  The characterisation used by
  \cite{LSZ2025,BV2025} --- $\zeta_p(s)=\lim\zeta(k)$ along negative integers
  $k\equiv s\pmod{p-1}$ with $k\to s$ $p$-adically, i.e.\ $-B_n/n$ at
  $n\equiv1-s\pmod{\lcm(p^N,p-1)}$ --- gives \VER{$16/16$ matches at
  $p\in\{5,7,11,13\}$, $s\in\{3,5\}$, $N\in\{2,3\}$; e.g.\ $p=13$, $s=5$,
  $n=26360$}.
\end{itemize}

\subsection{Values}\label{app:KLvalues}

Base-$p$ digits, least significant first (Table~\ref{tab:zetadigits}).  These
are the inputs to Tables~\ref{tab:lll1} and \ref{tab:lll2}.

\begin{table}[h]
\centering\small
\caption{$\zeta_p(3)$, base-$p$ digits, least significant first.}
\label{tab:zetadigits}
\begin{tabular}{@{}cl@{}}
\toprule
$p$ & digits of $\zeta_p(3)$\\
\midrule
5 & 2 4 1 2 3 1 4 2 0 2 0 0 2 3 3 0 $\dots$\\
7 & 1 3 5 4 1 0 2 0 2 4 1 2 5 0 4 6 $\dots$\\
11 & 5 4 1 0 5 8 5 2 6 8 3 1 4 3\\
13 & 6 6 5 2 11 12 5 1 5 2 7 5 10 10\\
\bottomrule
\end{tabular}
\end{table}

\section{The Lai--Sprang--Zudilin denominator ledger, reproduced}\label{app:LSZ}

Problem~\ref{prob:bridge} rests on the assertion that the denominator budget of
the known $p$-adic construction is tight.  We verified this by regenerating the
construction of \cite{LSZ2025} exactly.

Their object is $R_n(t)=2^{8n}(2t+n)(t+1/2)_n^4/(t)_{n+1}^4$, of degree $-3$,
with $S_n=-\int_{\ZZ_2}R_n'(t+1/2)\,dt$; the two vanishing mechanisms recalled in
\S\ref{ss:consistency} reduce $S_n$ to a two-term form
$S_n=\rho_{n,0}+\rho_{n,3}\zeta_2(5)$.  The integrality of $\rho_{n,3}$ comes
from a very-well-poised ${}^9V_8$ differentiated once in $\varepsilon$, turned by
Andrews' transformation \cite{KrattenthalerRivoal2007} into the explicit
double sum
\[
  \rho_n=\!\!\sum_{0\le i\le k\le n}\!\!
    2^{4(n-k)}\bin{2i}{i}^2\bin{2n-2i}{n-i}\bin{2k-2i}{k-i}\bin{2k}{k}^2\bin{2n-2k}{n-k},
  \qquad \rho_{n,3}=768\rho_n,
\]
while the bound $\vp(\rho_{n,0})\ge-5$ comes from the same machine at
${}^{13}V_{12}$ with three $\varepsilon$-derivatives --- this is where the
``$-5$'' originates.  The smallness exponent $\alpha=16\log2$ comes from
Sprang's $\Delta$-operator \cite{Sprang2020} in the estimate
$v_2(\int f)\ge\Delta(f)-1$ of \cite[Lemma~2.4]{Lai2025},
and the growth exponent $\beta=8\log2+5$ from the
double root $2^8$ of the characteristic polynomial $\lambda^2-2^9\lambda+2^{16}$
together with the denominator cost $\Phi_n^{-1}d_n^6=e^{5n+o(n)}$.

\begin{finding}\label{find:LSZledger}
\VER{$n\le1000$, exact rational arithmetic}  Regenerating $\rho_{n,0}$,
$\rho_{n,3}$, $\rho_n$ from the recursion of \cite{LSZ2025}:
\begin{enumerate}[label=(\roman*),leftmargin=2em]
\item $\rho_{n,3}=768\rho_n$ and $\rho_n\in\ZZ$;
\item the exact Casoratian
  $\rho_{n,0}\rho_{n+1,3}-\rho_{n+1,0}\rho_{n,3}=3\cdot2^{16n+18}/(n+1)^5\ne0$,
  \VER{$n\le60$ by regenerating both sequences from the recursion};
\item the conjectured denominator condition $d_n^5\rho_{n,0}\in\ZZ$ holds, and
  $d_n^4$ does \emph{not} suffice at any tested $n$;
\item the per-prime ledger $e_p(n):=-\vp(\rho_{n,0})$ equals
  $5\floor{\log_pn}$ exactly for almost every $p$, with isolated deficits of
  $1$--$3$ at $p=3,13$;
\item $\log(\operatorname{den}\rho_{n,0})/n\to5$ from below, the gap shrinking
  from $0.138$ at $n=200$ to $0.042$ at $n=1000$.
\end{enumerate}
Consequently the $d_n^5$ budget is tight: there is no free denominator saving
hidden in the construction.
\end{finding}

\subsection{The margin ledger}\label{app:margin}

The three published individual $p$-adic irrationality results have the same
shape --- smallness $\alpha=2\cdot(\text{growth})$, denominator cost equal to
the weight $w$, and $\text{growth}_p=(w+3)\log p/(p-1)$ --- giving
\[
  \mathrm{margin}(p,w)=\alpha-\beta=(w+3)\frac{\log p}{p-1}-w,
  \qquad \mu\le\frac{\alpha}{\alpha-\beta}.
\]

\begin{finding}\label{find:ledgerfit}
\VER{the fit reproduces all three published irrationality measures exactly}
See Table~\ref{tab:ledger}.
\end{finding}

\begin{table}[h]
\centering\small
\caption{The ledger fit against the published measures.}
\label{tab:ledger}
\begin{tabular}{@{}lllll@{}}
\toprule
known result & $\alpha$ & $\beta$ & $\mu=\alpha/(\alpha-\beta)$ & published\\
\midrule
$\zeta_2(3)$ \cite{Beukers2008,Calegari2005} & $12\log2$ & $6\log2+3$ & $7.1773988991\dots$ & $7.177398\dots$\\
$\zeta_3(3)$ \cite{Beukers2008,Calegari2005} & $6\log3$ & $3\log3+3$ & $22.2814479514\dots$ & $22.281447\dots$\\
$\zeta_2(5)$ \cite{LSZ2025} & $16\log2$ & $8\log2+5$ & $20.3426517389\dots$ & $20.342651\dots$\\
\bottomrule
\end{tabular}
\end{table}

Extrapolating the same ledger (\FIT: three data points, all at $p\in\{2,3\}$)
gives the margins of Table~\ref{tab:margins}.  The \emph{conclusion} is robust
independently of the fit, in that nothing individual is known for any $p\ge5$.

\begin{table}[h]
\centering\small
\caption{Extrapolated margins $(w+3)\log p/(p-1)-w$.  Positive means a proof
exists by this route.}
\label{tab:margins}
\begin{tabular}{@{}ccc@{}}
\toprule
$p$ & margin, weight 3 & margin, weight 5\\
\midrule
2 & $+1.1589$ (known) & $+0.5452$ (known)\\
3 & $+0.2958$ (known) & $-0.6056$\\
5 & $-0.5858$ & $-1.7811$\\
7 & $-1.0541$ & $-2.4055$\\
11 & $-1.5613$ & ---\\
13 & $-1.7175$ & ---\\
\bottomrule
\end{tabular}
\end{table}

\section{Numerical data}\label{app:data}

Table~\ref{tab:constants} records the tower limits and the derived constants in
base-$p$ digits, least significant first, computed from exact rationals.
Certified precisions: Ap\'ery $24$ digits at $p=5$ (from $n=78\,125$; $27$ from
$n=400\,000$), $21$ at $p=7$, $15$--$18$ at $p=11,13$; Brown--Zudilin $18$ at
$p=5$ (from $n=3125$), $15$ at $p=7$ (from $n=2401$), $12$ at $p=11,13$.

The Ap\'ery column of Table~\ref{tab:constants} was computed twice, by
independent programs using different algorithms (an $O(N^2)$ exact-rational
recurrence and an $O(N)$ mod-$p^K$ pass), with identical results.

\begin{table}[h]
\centering\scriptsize
\setlength{\tabcolsep}{3pt}
\caption{Base-$p$ digits, least significant first, $12$ digits shown; every
constant is normalised by $p^{-v}$ where $v$ is its valuation.  Entries marked
$\dagger$ have $v=-1$; all others have $v=0$.}
\label{tab:constants}
\begin{tabular}{@{}cllll@{}}
\toprule
$p$ & $\Lambda_1$ (Ap\'ery) & $\Lap_1$ (BZ) & $\Lah_1$ (BZ) & $\widehat{c}_p(1)=\Lah_1/\Lap_1$\\
\midrule
5 & 1 1 3 3 4 0 1 1 1 0 2 2$^\dagger$ & 3 3 2 2 4 0 0 2 3 0 1 3 & 4 2 3 0 1 4 4 3 3 4 1 1 & 3 4 2 0 4 3 0 4 0 4 2 2\\
7 & 4 1 4 4 6 3 0 5 2 6 4 5 & 2 6 5 1 3 3 2 4 0 6 5 1$^\dagger$ & 2 5 4 5 1 6 3 6 2 0 1 0$^\dagger$ & 1 3 4 5 4 4 2 6 2 5 4 1\\
11 & 10 8 8 4 10 7 0 6 5 5 5 4 & 3 1 5 7 4 7 2 9 7 9 1 10 & 5 10 6 5 7 3 8 2 3 0 0 0 & 9 7 6 9 1 7 7 5 9 4 9 5\\
13 & 9 2 5 9 6 6 7 1 10 8 11 7 & 8 12 6 7 8 3 12 2 12 10 12 3 & 6 10 4 2 0 7 7 8 8 11 0 11 & 4 8 4 6 11 8 11 11 10 5 11 10\\
\bottomrule
\end{tabular}
\end{table}

\begin{table}[h]
\centering\scriptsize
\setlength{\tabcolsep}{3pt}
\caption{The ingredients of the structure theorem, base-$p$ digits, least
significant first, each normalised by $p^{-v}$.  Entries marked $\dagger$ have
$v=-1$, entries marked $\ddagger$ have $v=+1$; all others have $v=0$.}
\label{tab:ingredients}
\begin{tabular}{@{}cllll@{}}
\toprule
$p$ & $\Lambda_1$ & $\Atil_1$ & $h_p(1)$ & $c_p(2,1)$\\
\midrule
5 & 1 1 3 3 4 0 1 1 1 0 2 2$^\dagger$ & 1 0 2 0 2 4 3 0 1 0 3 1$^\ddagger$ & 1 0 0 0 4 3 1 0 3 2 1 0 & 2 0 0 2 0 0 2 2 2 2 3 1\\
7 & 4 1 4 4 6 3 0 5 2 6 4 5 & 5 0 0 0 5 3 3 0 4 1 1 6 & 1 0 0 0 0 4 3 6 1 3 5 4 & 2 0 0 3 1 0 3 1 6 5 2 3\\
11 & 10 8 8 4 10 7 0 6 5 5 5 4 & 5 0 0 7 3 5 6 0 3 9 8 3 & 1 0 0 0 0 5 5 10 3 0 10 0 & 2 0 0 2 4 4 2 4 5 3 2 0\\
13 & 9 2 5 9 6 6 7 1 10 8 11 7 & 5 0 0 7 0 4 11 12 4 6 3 8 & 1 0 0 0 0 12 11 4 9 12 8 4 & 2 0 0 2 0 2 4 0 4 7 5 8\\
\bottomrule
\end{tabular}
\end{table}

\section{Reproduction}\label{app:repro}

All computations use exact integer or rational arithmetic; no floating point
occurs at any stage, and every $p$-adic digit reported above is derived from an
exact rational.  The programs are organised as follows.

\begin{itemize}[leftmargin=1.5em]
\item \emph{Sequence generation}: exact Ap\'ery and sporadic pairs from the
  integer recurrence \eqref{eq:int-rec} and its analogues; exact extension of
  the Brown--Zudilin ladders $Q,P,\Ph$ by the certified order-$3$ recurrence,
  with the recurrence residual checked to be exactly $0$ at every $n\le357$ on
  every row.
\item \emph{$p$-adic arithmetic}: a $\Qp$ class carrying an honest precision
  bound through every operation, so that no digit is ever reported that is not
  certified.
\item \emph{Tower data}: $L_s$, $\Lambda_a$, $\Atil_a$, $\Btil_a$ for all towers
  and branches, with the per-level agreement recorded (Tables~\ref{tab:aperyrate},
  \ref{tab:bzrate}).
\item \emph{Kubota--Leopoldt}: \eqref{eq:KL} plus both cross-checks of
  Appendix~\ref{app:KLvalid}.
\item \emph{Integer relations}: an exact-Fraction LLL with weighted residue
  column, plus rational reconstruction, self-tested on planted relations before
  use.
\item \emph{Identification}: $h_p$ from Theorem~\ref{thm:hp}; $c_p(A,B)$ both
  directly (as a limit) and from \eqref{eq:kaz}; the block-by-block check of
  Table~\ref{tab:blocks}.
\end{itemize}

\subsection*{Acknowledgements}

[TO BE SUPPLIED]

% =====================================================================
\begin{thebibliography}{99}
% =====================================================================

% --- All entries below have been checked against a fetched source (arXiv
%     metadata, publisher record, or the text of the cited work); see the
%     verification memo work/BIBLIO_VERIFY.md for the evidence trail.

\bibitem{PaperA-inprep}
[AUTHOR],
\emph{Digit-descent congruences for Ap\'ery-like pairs: the twisted law
$(\mathrm{LB}_w^\chi)$}.
In preparation.  (Paper~A of this series.)

\bibitem{PaperB-inprep}
[AUTHOR],
\emph{The purity defect of the Brown--Zudilin $\zeta(5)$ family}.
In preparation.  (Paper~B of this series.)

\bibitem{ASvSZ2008}
G.~Almkvist, D.~van Straten and W.~Zudilin,
\emph{Ap\'ery limits of differential equations of order $4$ and $5$},
in: Modular Forms and String Duality, Fields Inst.\ Commun.\ \textbf{54},
Amer.\ Math.\ Soc., Providence, RI, 2008, pp.~105--123.

\bibitem{Apery1979}
R.~Ap\'ery,
\emph{Irrationalit\'e de $\zeta(2)$ et $\zeta(3)$},
in: Journ\'ees Arithm\'etiques (Luminy, 1978),
Ast\'erisque \textbf{61} (1979), 11--13.

\bibitem{Beukers1987}
F.~Beukers,
\emph{Another congruence for the Ap\'ery numbers},
J.\ Number Theory \textbf{25} (1987), no.~2, 201--210;
see also F.~Beukers, \emph{Irrationality proofs using modular forms},
Ast\'erisque \textbf{147--148} (1987), 271--283.

\bibitem{Beukers2008}
F.~Beukers,
\emph{Irrationality of some $p$-adic $L$-values},
Acta Math.\ Sin.\ (Engl.\ Ser.) \textbf{24} (2008), no.~4, 663--686.
(The measures for $\zeta_2(3)$, $\zeta_3(3)$ quoted in
Table~\ref{tab:ledger} are Theorem~11.2 there.)

\bibitem{BV2025}
F.~Beukers and M.~Vlasenko,
\emph{Frobenius structure and $p$-adic zeta values},
Adv.\ Math.\ \textbf{480} (2025); arXiv:2302.09603.

\bibitem{BVdwork3}
F.~Beukers and M.~Vlasenko,
\emph{Dwork crystals III: from excellent Frobenius lifts towards
supercongruences},
Int.\ Math.\ Res.\ Not.\ (2023); arXiv:2105.14841.

\bibitem{BrownZudilin2022}
F.~Brown and W.~Zudilin,
\emph{On cellular rational approximations to $\zeta(5)$},
arXiv:2210.03391.

\bibitem{Brown2014}
F.~Brown,
\emph{Irrationality proofs for zeta values, moduli spaces and dinner parties},
arXiv:1412.6508.

\bibitem{Calegari2005}
F.~Calegari,
\emph{Irrationality of certain $p$-adic periods for small $p$},
Int.\ Math.\ Res.\ Not.\ (2005), no.~20, 1235--1249.

\bibitem{ChamberlandStraub}
M.~Chamberland and A.~Straub,
\emph{Ap\'ery limits: experiments and proofs},
Amer.\ Math.\ Monthly \textbf{128} (2021), no.~9, 811--824;
arXiv:2011.03400.

\bibitem{Cooper2012}
S.~Cooper,
\emph{Sporadic sequences, modular forms and new series for $1/\pi$},
Ramanujan J.\ \textbf{29} (2012), no.~1--3, 163--183.
(The three additional sporadic solutions $s_7$, $s_{10}$, $s_{18}$ are
attributed to this paper by \cite[\S2]{MalikStraub}.)

\bibitem{Fischler2003}
S.~Fischler,
\emph{Irrationalit\'e de valeurs de z\^eta (d'apr\`es Ap\'ery, Rivoal, \dots)},
S\'eminaire Bourbaki exp.\ 910, Ast\'erisque \textbf{294} (2004), 27--62;
arXiv:math/0303066.

\bibitem{FSZ2018}
S.~Fischler, J.~Sprang and W.~Zudilin,
\emph{Many odd zeta values are irrational},
Compositio Math.\ \textbf{155} (2019), no.~5, 938--952; arXiv:1803.08905.

\bibitem{Gessel1982}
I.~M.~Gessel,
\emph{Some congruences for Ap\'ery numbers},
J.\ Number Theory \textbf{14} (1982), no.~3, 362--368.

\bibitem{GolyshevZagier}
V.~Golyshev and D.~Zagier,
\emph{Interpolated Ap\'ery numbers, quasiperiods of modular forms, and motivic
gamma functions},
in: Integrability, Quantization, and Geometry II,
Proc.\ Sympos.\ Pure Math.\ \textbf{103.2},
Amer.\ Math.\ Soc., 2021, pp.~281--301.

\bibitem{GrossKoblitz1979}
B.~H.~Gross and N.~Koblitz,
\emph{Gauss sums and the $p$-adic $\Gamma$-function},
Ann.\ of Math.\ (2) \textbf{109} (1979), no.~3, 569--581.

\bibitem{JacobsthalKazandzidis}
V.~Brun, J.~O.~Stubban, J.~E.~Fjeldstad, R.~Tambs Lyche, K.~E.~Aubert,
W.~Ljunggren and E.~Jacobsthal,
\emph{On the divisibility of the difference between two binomial coefficients},
Den 11te Skandinaviske Matematikerkongress (Trondheim, 1949),
Johan Grundt Tanums Forlag, Oslo, 1952, pp.~42--54;
G.~S.~Kazandzidis,
\emph{Congruences on the binomial coefficients},
Bull.\ Soc.\ Math.\ Gr\`ece (N.S.) \textbf{9} (fasc.~1) (1968), 1--12.

\bibitem{Kerr}
M.~Kerr,
\emph{Unipotent extensions and differential equations (after Bloch--Vlasenko)},
arXiv:2008.03618.

\bibitem{KrattenthalerRivoal2007}
C.~Krattenthaler and T.~Rivoal,
\emph{Hyperg\'eom\'etrie et fonction z\^eta de Riemann},
Mem.\ Amer.\ Math.\ Soc.\ \textbf{186} (2007), no.~875, x\,+\,87~pp.
(Andrews' transformation is Th\'eor\`eme~8 there; it is applied in this form
in \cite{LSZ2025}.)

\bibitem{KubotaLeopoldt1964}
T.~Kubota and H.~W.~Leopoldt,
\emph{Eine $p$-adische Theorie der Zetawerte. Teil~I: Einf\"uhrung der
$p$-adischen Dirichletschen $L$-Funktionen},
J.\ Reine Angew.\ Math.\ \textbf{214/215} (1964), 328--339.

\bibitem{Lai2025}
L.~Lai,
\emph{On the irrationality of certain $2$-adic zeta values},
Int.\ J. Number Theory \textbf{21} (2025), no.~1, 207--235.

\bibitem{LaiSprang2023}
L.~Lai and J.~Sprang,
\emph{Many $p$-adic odd zeta values are irrational},
arXiv:2306.10393.

\bibitem{LSZ2025}
L.~Lai, J.~Sprang and W.~Zudilin,
\emph{A note on the irrationality of $\zeta_2(5)$},
arXiv:2505.05005.

\bibitem{MalikStraub}
A.~Malik and A.~Straub,
\emph{Divisibility properties of sporadic Ap\'ery-like numbers},
arXiv:1508.00297.

\bibitem{Morita1975}
Y.~Morita,
\emph{A $p$-adic analogue of the $\Gamma$-function},
J.\ Fac.\ Sci.\ Univ.\ Tokyo Sect.\ IA Math.\ \textbf{22} (1975), no.~2,
255--266.
(The displayed product formula is the iterate of Morita's relation
$\Gp(n+1)=(-1)^{n+1}n!\,/\,(p^{\floor{n/p}}\floor{n/p}!)$.)

\bibitem{RhinViola}
G.~Rhin and C.~Viola,
\emph{On a permutation group related to $\zeta(2)$},
Acta Arith.\ \textbf{77} (1996), no.~1, 23--56;
\emph{The group structure for $\zeta(3)$},
Acta Arith.\ \textbf{97} (2001), no.~3, 269--293.

\bibitem{RoyVlasenko}
B.~Roy and M.~Vlasenko,
\emph{Frobenius constants for families of elliptic curves},
arXiv:2206.15181.

\bibitem{Sprang2020}
J.~Sprang,
\emph{Linear independence result for $p$-adic $L$-values},
Duke Math.\ J. \textbf{169} (2020), no.~18, 3439--3476.

\bibitem{Straub2023}
A.~Straub,
\emph{Gessel--Lucas congruences for sporadic sequences},
Monatsh.\ Math.\ (2023); arXiv:2301.12248.

\bibitem{Supercong}
F.~Beukers,
\emph{Some congruences for the Ap\'ery numbers},
J.\ Number Theory \textbf{21} (1985), no.~2, 141--155;
M.~J.~Coster,
\emph{Supercongruences},
Ph.D.\ thesis, Universiteit Leiden, 1988;
T.~Ishikawa,
\emph{Super congruence for the Ap\'ery numbers},
Nagoya Math.\ J. \textbf{118} (1990), 195--202.

\bibitem{Zagier2009}
D.~Zagier,
\emph{Integral solutions of Ap\'ery-like recurrence equations},
in: J.~Harnad and P.~Winternitz (eds.), Groups and Symmetries:
From Neolithic Scots to John McKay,
CRM Proc.\ Lecture Notes \textbf{47}, Amer.\ Math.\ Soc., 2009, pp.~349--366.

\end{thebibliography}

\end{document}
